 % 本文件由 scripts/assemble.py 自动装配，勿手改（改 parts/ 后重跑）
% 第10章 磁性

\chapter{磁性}

\epigraphbook{万物皆始于秩序，故亦将终于秩序，并将依循秩序的制定者、\\天上之城的神妙数学，再度开始。}{SIR THOMAS BROWNE}

\section{轨道磁化率}

在本章中，我们所关注的是固体的磁化率本身，而不涉及磁场对固体其他性质（例如电导率）的影响。这是一个很大的课题，我们无法细致处理。例如，我们将略去大多数磁共振（magnetic resonance）现象。我们也不讨论技术上十分重要的、可永久磁化的材料的那些宏观性质的理论，例如磁滞（hysteresis）、矫顽力（coercivity）、剩磁（remanent magnetism）等等；这些性质本质上取决于样品被自发磁化而分成若干磁畴（domains），各磁畴的磁矩取向彼此不同。我们甚至不打算考虑划分磁化不同的磁畴的 Bloch 壁（Bloch wall）的优美理论——这种壁在施加磁场时还可以被驱使移动。

我们首先考虑\emph{非磁性固体}，在其中完全不必怀疑原子或离子携带有局域磁矩。我们也忽略电子自旋，或者假定所有自旋都紧紧地配对。即便如此，电子的轨道运动仍会对磁化率有贡献。

最简单的情形是每个原子或离子都由电子的闭合壳层组成，且激发到更高状态的能量很大。众所周知，这就导致抗磁性（diamagnetism）；磁化率为负，近似地由下式给出：
\begin{equation*}
\chi=-\frac{Ze^{2}N}{6mc^{2}}\,\overline{r^{2}},
\tag{10.1}
\end{equation*}
其中 $Z$ 是原子中电子的总数，$N$ 是单位体积中的原子数，$\overline{r^{2}}$ 是每个原子周围电子电荷云的方均半径。这正是我们处理稀有气体固体、或者离子晶体中离子的贡献时应当使用的那类公式。

上述结果本质上来自一阶微扰能量
\begin{equation*}
\left\langle 0\,\left|\frac{e^{2}}{2mc^{2}}A^{2}\right|\,0\right\rangle,
\tag{10.2}
\end{equation*}
其中 $\mathbf{A}$ 是磁场的矢量势，$|0\rangle$ 是每个原子的基态。当我们按照磁场中哈密顿量的标准表达式，把动能写成
\begin{equation*}
T=\frac{1}{2m}\left(\frac{\hbar}{i}\nabla-\frac{e}{c}\mathbf{A}\right)^{2}
\tag{10.3}
\end{equation*}
的形式时（如式 (9.51) 那样），上式便自动出现。

但是，如果原子存在激发态，比如说 $|n\rangle$，其轨道矩不为零，那么 (10.3) 中线性于 $\mathbf{A}$ 的项——也就是给出算符
\begin{equation*}
\frac{e}{2mc}\,\mathbf{H}\cdot\mathbf{L}=\mathbf{H}\cdot\boldsymbol{\mu}_{L},
\tag{10.4}
\end{equation*}
的那一项——就可能把这些激发态混入基态，其中 $\mathbf{L}$ 是电子的轨道角动量算符，而 $\boldsymbol{\mu}_{L}$ 是与之成比例的通常的磁矩算符。在二阶微扰中，这一项对能量的贡献为
\begin{equation*}
\delta\mathcal{E}=\sum_{n}\frac{|\langle 0|\mathbf{H}\cdot\boldsymbol{\mu}_{L}|n\rangle|^{2}}{\mathcal{E}_{0}-\mathcal{E}_{n}}
\tag{10.5}
\end{equation*}
它是正的，并且是 $H$ 的二次式。这就是原子、离子或分子的恒定的 van Vleck 顺磁性（van Vleck paramagnetism）的来源。

实践表明，把 (10.2) 和 (10.5) 应用于实际固体极其复杂；这就是我们只在原则上引用这些公式的原因。但我们从 (10.5) 中注意到，当激发能量变小时——例如在填满的能带之上有小的带隙时——这一贡献应当增大。当这个带隙趋于零，也就是当我们拥有自由载流子时，会发生什么呢？

这时理论又变得非常复杂，但对自由电子却存在精确解。在 §9.6 中我们导出了磁场中自由电子的能级；只要贯彻那套导致 de Haas--van Alphen 效应的 §9.7 的形式框架，算出整个电子气能量的平均变化并不困难。我们不寻求自由能中的振荡项，而是集中于那些作为 $H$ 的单调函数的项。例如，在 (9.80) 中我们略去了
\begin{equation*}
\frac{1}{4\pi^{2}s^{2}kT}\left[f^{0}(\mathcal{E})\frac{\partial\mathcal{E}}{\partial x}\right]_{x=0}
=\frac{1}{4\pi^{2}s^{2}kT}\,\frac{e\hbar H}{mc}\left[f^{0}(\mathcal{E})\right]_{x=0}
\tag{10.6}
\end{equation*}
（其中利用了 (9.82) 和 (9.58)）。

但这是 (9.79) 中的一项，而 (9.79) 中含有对所有 $s$ 值的求和以及对 Fermi 球的所有切片 $k_{z}$ 的积分。于是，对自由能的总贡献（除去 (9.79) 中与 $H$ 无关的前两项）便是
\begin{align*}
\Delta F_{L} &=-2kT\sum_{s=1}^{\infty}(-1)^{s}\int_{-\infty}^{\infty}\frac{eH}{2\pi^{2}\hbar c}\,\frac{1}{4\pi^{2}s^{2}kT}\,\frac{e\hbar H}{mc}\left[f^{0}(\mathcal{E})\right]_{x=0}dk_{z}\\
&=-\frac{1}{4\pi^{4}}\left(\frac{eH}{c}\right)^{2}\frac{1}{m}\int_{-\infty}^{\infty}\left[f^{0}(\mathcal{E})\right]_{x=0}dk_{z}\sum_{s=1}^{\infty}\frac{(-1)^{s}}{s^{2}}.
\tag{10.7}
\end{align*}
对 $k_{z}$ 的积分本质上量度的是 k 空间中沿 Fermi 球轴线的填充区域的长度，而对 $s$ 的求和给出 $-\frac{1}{12}\pi^{2}$。把它表示成磁化率，我们便得到 \emph{Landau 抗磁性}（Landau diamagnetism）
\begin{equation*}
\chi_{L}=-\frac{e^{2}k_{F}}{12\pi^{2}mc^{2}}.
\tag{10.8}
\end{equation*}

这个结果为负，是因为能级的聚束倾向于增大系统的总能量。从图 171 可以看到这一点。只有当 费米能级恰好位于两个量子化能级的正中间时，电子才能够“凝聚”到各自相应的管道上而不改变其平均能量。当一个量子化能级逼近 费米能级时，它会把电子一同带上去；这就是 $\Delta F_{L}$ 为正的原因。磁化率与温度无关，因为这是一种平均效应，并不要求能级间隔 $\hbar\omega_{H}$ 大于 $kT$。

按照 §9.7 论证的思路，对任意 费米面这一更一般的情形来计算 (10.7) 并不困难。但这样得到的结果并不正确。周期晶格中电子的 Landau 抗磁性需要更复杂的分析，我们在此不拟给出。此外，在 van Vleck 贡献与 Landau 贡献之间还存在交叉项。

\section{自旋顺磁性}

共享同一轨道态的、自旋相反的电子之间的简并，会被磁场解除。在金属中，这引起电子在两种自旋取向之间重新分布，从而产生一个磁矩。这个效应的计算是初等的。假设我们区分 + 自旋与 $-$ 自旋的电子（相对于 $\mathbf{H}$ 的方向）。那么它们的能量将为
\begin{equation*}
\left.\begin{aligned}
\mathcal{E}_{\mathbf{k}+} &= \mathcal{E}(\mathbf{k})-\mu_{0}H,\\
\mathcal{E}_{\mathbf{k}-} &= \mathcal{E}(\mathbf{k})+\mu_{0}H,
\end{aligned}\right\}
\tag{10.9}
\end{equation*}
其中 $\mathcal{E}(\mathbf{k})$ 是不存在磁场时的能量，而 $\mu_{0}$ 是电子的磁矩。

每个状态中的电子数将由两个不同的 Fermi--Dirac 分布给出，如同式 (4.8) 那样，但具有相同的化学势 $\zeta$。于是，正如我们在图 175 中看到的，+ 自旋态容纳的电子比先前要多——总数为
\begin{equation*}
n_{+}=\int_{0}^{\infty}\frac{1}{2}\,\mathcal{N}(\mathcal{E})\,f^{0}(\mathcal{E}_{\mathbf{k}+})\,d\mathcal{E},
\tag{10.10}
\end{equation*}
这是因为在原来能量 $\mathcal{E}$ 处的态密度，是在 + 自旋与 $-$ 自旋的电子之间平分的。

\begin{figure}[H]
\centering
\includegraphics[width=0.72\textwidth]{fig_175.pdf}
\caption*{图 175\quad 磁场中“上”自旋与“下”自旋的 费米分布。}
\end{figure}

对于负自旋电子的密度有相应的公式。两者之差给出一个磁矩
\begin{align*}
M &=\mu_{0}(n_{+}-n_{-})\\
&=\mu_{0}\int_{0}^{\infty}\tfrac{1}{2}\{f_{0}(\mathcal{E}-\mu_{0}H)-f_{0}(\mathcal{E}+\mu_{0}H)\}\,\mathcal{N}(\mathcal{E})\,d\mathcal{E}\\
&\approx\mu_{0}^{2}H\int_{0}^{\infty}\left(-\frac{\partial f^{0}}{\partial\mathcal{E}}\right)\mathcal{N}(\mathcal{E})\,d\mathcal{E}\\
&=\mu_{0}^{2}H\,\mathcal{N}(\mathcal{E}_{F})
\tag{10.11}
\end{align*}
由式 (4.18)。

这个正的磁矩导致 Pauli 顺磁性（Pauli paramagnetism）。它不依赖于温度，并且直接量度 费米能级处的态密度。因此，原则上它应当与电子比热（§4.7）直接相关联；不过，当人们对多电子相互作用作出修正之后（§5.8），这种严格的等价性便不复存在。在实践中做这类比较的困难在于，抗磁性修正 (10.8) 并非可以忽略，而人们对它又知道得不够准确。不过在有些情形下，可以通过观察传导电子的电子自旋共振，直接测量自旋磁化率。

对自由电子气，容易证明
\begin{align*}
\chi_{P} &=\mu_{0}^{2}\,\mathcal{N}(\mathcal{E}_{F})\\
&=\frac{e^{2}k_{F}}{4\pi^{2}mc^{2}},
\tag{10.12}
\end{align*}
它恰好是 Landau 项 (10.8) 数值的三倍。但是，晶格对这两类磁化率的影响却截然不同。例如，假设 费米面是球形的，但“有效质量”变为 $m^{*}$。容易证明，$\chi_{P}$ 依赖于态密度，因而正比于 $m^{*}$；而 $\chi_{L}$ 由电子的轨道运动决定，不经过 Bohr 磁子的介入，所以将\emph{反比}于 $m^{*}$。

与自旋磁化率相联系的另一个现象是 Knight 移动（Knight shift）。在磁场 $\mathbf{H}$ 中，电子自旋被极化，每个电子具有平均磁矩
\begin{equation*}
\langle\boldsymbol{\mu}_{e}\rangle=\frac{1}{n}\,\chi_{P}\,\mathbf{H}.
\tag{10.13}
\end{equation*}
位于 $\mathbf{r}$ 处的一个电子，将与位于 $\mathbf{R}$ 处的原子核的磁矩 $\boldsymbol{\mu}_{I}$ 通过“接触相互作用”（contact interaction）
\begin{equation*}
\mathcal{H}_{\text{int.}}=\tfrac{8}{3}\pi\,(\boldsymbol{\mu}_{I}\cdot\boldsymbol{\mu}_{e})\,\delta(\mathbf{r}-\mathbf{R}),
\tag{10.14}
\end{equation*}
发生作用，这个效应在核共振理论中是众所周知的。把 (10.13) 代入 (10.14)，我们得到一个平均能量
\begin{align*}
\langle\mathcal{H}_{\text{int.}}\rangle &=\tfrac{8}{3}\pi\,\frac{\chi_{P}}{n}\,H\,|\psi(0)|^{2}\,\mu_{I}\\
&=\Delta\mu_{I}\cdot\mathbf{H},
\tag{10.15}
\end{align*}
姑且这样记，其中 $|\psi(0)|^{2}$ 是波函数在核处的模方。这个能量被解释为核的表观磁矩出现一个改变量 $\Delta\mu_{I}$——即它在磁场中共振频率的移动。

显然，Knight 移动应当正比于 $\chi_{P}$。但还有一个因子 $|\psi(0)|^{2}$，它无法由任何独立的实验给出。原则上这只涉及 费米面上的态，因此人们说它“量度了 费米能级处波函数中 $s$ 型成分的多少”。Knight 移动是一个容易测定的好物理量，但它的结果却不那么容易解释。

\section{居里--外斯定律与铁磁性}

现在我们假设每个原子的行为像一个小磁体，磁矩为 $\boldsymbol{\mu}$，例如可以来自过渡金属或稀土元素的离子中 $d$ 电子或 $f$ 电子的未满壳层。在磁场 $\mathbf{H}$ 中，每个小磁体将具有能量 $-\boldsymbol{\mu}\cdot\mathbf{H}$。假设每个原子都与近邻无关；由于它是局域的，它满足 Boltzmann 统计。磁矩为 $\boldsymbol{\mu}$ 的原子所占的比例将正比于
\begin{equation*}
n(\boldsymbol{\mu})=e^{\boldsymbol{\mu}\cdot\mathbf{H}/kT}.
\tag{10.16}
\end{equation*}

如果我们假设每个磁体都可以自由转动，那么平均磁矩必定是
\begin{equation*}
\langle\boldsymbol{\mu}\rangle=\int\boldsymbol{\mu}\,e^{\boldsymbol{\mu}\cdot\mathbf{H}/kT}d\Omega\,\Big/\int e^{\boldsymbol{\mu}\cdot\mathbf{H}/kT}d\Omega,
\tag{10.17}
\end{equation*}
其中 $d\Omega$ 是 $\boldsymbol{\mu}$ 转动时的立体角元。单位体积内 $N$ 个这种原子的集合将具有磁化率
\begin{align*}
\chi &=N\frac{\partial\langle\boldsymbol{\mu}\rangle}{\partial\mathbf{H}}=\int(N/kT)\,\boldsymbol{\mu}\boldsymbol{\mu}\,e^{\boldsymbol{\mu}\cdot\mathbf{H}/kT}d\Omega\,\Big/\int e^{\boldsymbol{\mu}\cdot\mathbf{H}/kT}d\Omega\\
&\approx\frac{N}{kT}\langle\boldsymbol{\mu}\boldsymbol{\mu}\rangle=\frac{1}{3}\,\frac{N\langle\mu^{2}\rangle}{kT}
\tag{10.18}
\end{align*}
此时 $H$ 已小到指数项可以忽略；并且我们已经把并矢张量 $\boldsymbol{\mu}\boldsymbol{\mu}$ 的平均，同矢量 $\boldsymbol{\mu}$ 之长度的平方 $\mu^{2}$ 的平均作了比较。

如今对我们来说，更自然的做法是把每个磁体同一个自旋角动量 $\mathbf{S}$ 联系起来，它沿磁场方向量子化：
\begin{equation*}
\boldsymbol{\mu}\cdot\mathbf{H}=\beta HS_{z},
\tag{10.19}
\end{equation*}
其中 $S_{z}$ 以整数台阶从 $S$ 变到 $-S$，而 $\beta$ 是 Bohr 磁子的两倍。例如，设 $S=\frac{1}{2}$。这时在 (10.17) 中我们遇到的是求和而不是积分：
\begin{align*}
\langle\boldsymbol{\mu}\rangle &=\beta S\left[\frac{e^{\beta SH/kT}-e^{-\beta SH/kT}}{e^{\beta SH/kT}+e^{-\beta SH/kT}}\right]\\
&=\beta S\,\tanh\frac{\beta SH}{kT}.
\tag{10.20}
\end{align*}

在 $H$ 很小的极限下，这个结果便与 (10.18) 相同，只要我们注意到，按照总角动量大小的规则，
\begin{equation*}
\langle\mu^{2}\rangle=S(S+1)\beta^{2}.
\tag{10.21}
\end{equation*}

结果 (10.18) 就是著名的顺磁磁化率的 \emph{Curie 定律}（Curie Law）。某些类型的固体近似地服从它，尤其是磁性原子或离子在晶体中彼此相距很远的那些固体，因为这时独立性假设成立。

然而在许多物质中，近邻自旋之间存在相互作用。为了近似地把这一点考虑进去，可以引入 \emph{Weiss 场}（Weiss field）
\begin{equation*}
\mathbf{H}_{I}=\lambda N\langle\boldsymbol{\mu}\rangle.
\tag{10.22}
\end{equation*}
这个场被认为来自每个原子所处的磁化平均环境——不过，用来量度相互作用强度的参量 $\lambda$，我们让它保持任意。

作用在一个原子上的总场现在是 $\mathbf{H}+\mathbf{H}_{I}$。把它代替 $\mathbf{H}$ 代入 (10.20)，我们近似地得到
\begin{equation*}
N\langle\boldsymbol{\mu}\rangle=\frac{1}{3}\,\frac{N\langle\mu^{2}\rangle}{kT}\left(\mathbf{H}+\lambda N\langle\boldsymbol{\mu}\rangle\right),
\tag{10.23}
\end{equation*}
由此得
\begin{align*}
\chi &\equiv N\left\langle\frac{\partial\boldsymbol{\mu}}{\partial\mathbf{H}}\right\rangle\\
&=\frac{N\langle\mu^{2}\rangle}{3k(T-\theta)},
\tag{10.24}
\end{align*}
其中
\begin{equation*}
\theta=\frac{\lambda N\langle\mu^{2}\rangle}{3k}.
\tag{10.25}
\end{equation*}
这就是所谓的\emph{居里--外斯定律}（Curie--Weiss law），它描述许多固体的磁化率行为。

参量 $\theta$ 显然具有温度的量纲。如果在 (10.24) 中 $T<\theta$ 会怎么样？磁化率似乎会先变为无穷大，然后变为负值。我们必须回到 (10.20)，求解方程
\begin{equation*}
N\langle\boldsymbol{\mu}\rangle=N\beta S\,\tanh\left\{\frac{\beta S}{kT}\left(\lambda N\langle\boldsymbol{\mu}\rangle+H\right)\right\}.
\tag{10.26}
\end{equation*}
当 $H=0$ 时，通常的解是 $\langle\boldsymbol{\mu}\rangle=0$；但当 $T<\theta$ 时，我们发现还有另一个 $\langle\boldsymbol{\mu}\rangle\neq 0$ 的根。这个根很容易用图解法定位；在温度 $T=0$ 时它显然趋于
\begin{equation*}
\mathbf{M}(0)=N\beta S.
\tag{10.27}
\end{equation*}

\begin{figure}[H]
\centering
\includegraphics[width=0.72\textwidth]{fig_176.pdf}
\caption*{图 176\quad 铁磁体中的自发磁化。}
\end{figure}

这时系统的行为就好像每一个磁矩都平行排列一样；我们说这固体是\emph{铁磁性的}（ferromagnetic），并且看到，即使在不存在磁场时它也表现出很大的自发磁化。当我们使晶体升温时，自发磁化不断减小，直到在一个完全确定的温度——\emph{居里温度}（Curie temperature）——处消失：
\begin{equation*}
T_{\mathrm{C}}=\theta.
\tag{10.28}
\end{equation*}

\section{交换相互作用}

用居里温度 (10.25) 来衡量，Weiss 内场非常之大——远大于一个磁体阵列的普通内禀偶极场。在许多铁磁物质中，系数 $\lambda$ 并非 $\frac{4}{3}\pi$（对 Lorentz 场 (8.52) 而言它本可以取这个值），而是 1000 的量级。

通常把这个表观场归因于相邻原子上电子自旋之间的\emph{交换相互作用}（exchange interaction）。对于 Heisenberg 和 Dirac 的严格局域模型，我们可以给第 $\boldsymbol{l}$ 个格点上的原子指派一个总自旋算符 $\boldsymbol{\mathsf{S}}_{\boldsymbol{l}}$，并写下哈密顿量
\begin{equation*}
\mathcal{H}=-\sum_{\boldsymbol{l}\boldsymbol{l}'}J_{\boldsymbol{l}\boldsymbol{l}'}\,\boldsymbol{\mathsf{S}}_{\boldsymbol{l}}\cdot\boldsymbol{\mathsf{S}}_{\boldsymbol{l}'}-\beta\mathbf{H}\cdot\sum_{\boldsymbol{l}}\boldsymbol{\mathsf{S}}_{\boldsymbol{l}}.
\tag{10.29}
\end{equation*}
\emph{交换积分}（exchange integral）$J_{\boldsymbol{l}\boldsymbol{l}'}$ 是格点 $\boldsymbol{l}$ 与 $\boldsymbol{l}'$ 相对位置的函数，但只有当 $\boldsymbol{l}-\boldsymbol{l}'$ 为一两个格点间距时它才比较大。

我们可以把这个哈密顿量写成
\begin{equation*}
\mathcal{H}=-\sum_{\boldsymbol{l}}\bigl\{\sum_{\boldsymbol{l}'}J_{\boldsymbol{l}\boldsymbol{l}'}\boldsymbol{\mathsf{S}}_{\boldsymbol{l}'}+\beta\mathbf{H}\bigr\}\cdot\boldsymbol{\mathsf{S}}_{\boldsymbol{l}}.
\tag{10.30}
\end{equation*}
试想温度极低时会出现什么情形：那时所有自旋都排列整齐，以致对一切 $\boldsymbol{l}'$ 有
\begin{equation*}
\langle\beta\boldsymbol{\mathsf{S}}_{\boldsymbol{l}'}\rangle\approx\langle\boldsymbol{\mu}\rangle,
\tag{10.31}
\end{equation*}
于是我们可以把 (10.30) 中的第一项辨认成内场
\begin{align*}
\beta\mathbf{H}_{I}&=\sum_{\boldsymbol{l}'}J_{\boldsymbol{l}\boldsymbol{l}'}\boldsymbol{\mathsf{S}}_{\boldsymbol{l}'}\\
&=\frac{1}{\beta}\Bigl(\sum_{\boldsymbol{l}'}J_{\boldsymbol{l}\boldsymbol{l}'}\Bigr)\langle\boldsymbol{\mu}\rangle,
\tag{10.32}
\end{align*}
即
\begin{equation*}
\lambda=\frac{1}{N\beta^{2}}\Bigl(\sum_{\boldsymbol{l}'}J_{\boldsymbol{l}\boldsymbol{l}'}\Bigr)
\tag{10.33}
\end{equation*}
或者，用居里温度 (10.25) 来表示，
\begin{equation*}
k\theta=\frac{1}{3}S(S+1)\sum_{\boldsymbol{l}'}J_{\boldsymbol{l}\boldsymbol{l}'}.
\tag{10.34}
\end{equation*}
这便用系统转变为铁磁态所处的温度给出了交换相互作用的大小。

这个\emph{自旋哈密顿量}（spin Hamiltonian）(10.29) 的性质将是本章其余部分的主要论题。但是，它作为一个基本假设是否有充分的根据呢？对这个问题已有过许多不同的回答，而它正是铁磁性作为一种物理现象之谜的核心所在。

从群论的论证容易证明：相互作用取自旋算符的标量积这一形式，至少是最简单的可能形式，尽管还可能存在更复杂的其他项，例如\emph{偶极–偶极相互作用}（dipole–dipole interaction），
\begin{equation*}
\mathcal{H}_{\mathrm{dip.-dip.}}=\beta^{2}\biggl\{\frac{\boldsymbol{\mathsf{S}}_{1}\cdot\boldsymbol{\mathsf{S}}_{2}}{R^{3}}-\frac{3(\boldsymbol{\mathsf{S}}_{1}\cdot\mathbf{R})(\boldsymbol{\mathsf{S}}_{2}\cdot\mathbf{R})}{R^{5}}\biggr\},
\tag{10.35}
\end{equation*}
其中两个自旋相距 $\mathbf{R}$（一个矢量）。

同样众所周知的是，同一原子上或相邻原子上的电子自旋倾向于通过交换效应耦合起来——这是 Pauli 原理的一个推论。例如，若 $\phi_{a}$ 和 $\phi_{b}$ 是我们“放入两个电子”的两个波函数，那么按照这两个电子的自旋是平行还是反平行，我们可以构造出两类状态。这些状态就是对称与反对称组合
\begin{align*}
\Phi_{S}(\mathbf{r}_{1},\mathbf{r}_{2})&=2^{-\frac{1}{2}}\bigl\{\phi_{a}(\mathbf{r}_{1})\phi_{b}(\mathbf{r}_{2})+\phi_{a}(\mathbf{r}_{2})\phi_{b}(\mathbf{r}_{1})\bigr\},\\
\Phi_{A}(\mathbf{r}_{1},\mathbf{r}_{2})&=2^{-\frac{1}{2}}\bigl\{\phi_{a}(\mathbf{r}_{1})\phi_{b}(\mathbf{r}_{2})-\phi_{a}(\mathbf{r}_{2})\phi_{b}(\mathbf{r}_{1})\bigr\},
\tag{10.36}
\end{align*}
其中 $\Phi_{S}$ 属于反平行自旋（\emph{单态}，singlet state），而 $\Phi_{A}$ 属于平行自旋（\emph{三重态}，triplet state）。现在，如果我们计算库仑能 $e^{2}/|\mathbf{r}_{1}-\mathbf{r}_{2}|$ 在这两个态中的平均值，就会发现二者相差一个量
\begin{equation*}
2\iint\phi_{a}^{*}(\mathbf{r}_{1})\phi_{b}^{*}(\mathbf{r}_{2})\frac{e^{2}}{|\mathbf{r}_{1}-\mathbf{r}_{2}|}\phi_{a}(\mathbf{r}_{2})\phi_{b}(\mathbf{r}_{1})\,\mathrm{d}\mathbf{r}_{1}\,\mathrm{d}\mathbf{r}_{2},
\tag{10.37}
\end{equation*}
这就是交换积分。容易把这个差改写成如下形式，使它表示 $2J\,\boldsymbol{\mathsf{S}}_{1}\cdot\boldsymbol{\mathsf{S}}_{2}$ 在 $\boldsymbol{\mathsf{S}}_{1}$ 与 $\boldsymbol{\mathsf{S}}_{2}$ 平行和反平行两种情形下取值之差。于是，$J$ 的正负与大小便取决于这个积分的正负与大小。

有两种众所周知的情形。设两个电子属于同一原子，但并未填满一个闭合壳层。这时，库仑相互作用与原子轨道 $\phi_{a}$、$\phi_{b}$ 的形式使得 $J$ 为正；电子自旋倾向于排列起来，在与壳层中待填充的独立态数目相容的前提下使总自旋最大。这就是 Hund 规则（Hund's rule），它解释了为什么过渡金属离子的未满 d 壳层中的电子倾向于组合起来，赋予离子一个很大的永久磁矩（§5.5）。研究这类离子处于复杂晶体中各种环境时的顺磁性与磁共振性质，是物理学的一个重要分支。

但我们在这里关心的是不同离子上的电子自旋之间的相互作用——结果发现，在这种情形下算出的 $J$ 几乎总是负的，即倾向于使相邻格点上的自旋反平行。这方面最简单的例子是氢分子的 Heitler–London 模型，其中成键态的电子自旋是配对的。在 §4.2 中处理共价键时，我们曾把这一点视为理所当然。

因此，在这个基础上很难解释许多金属——通常是“过渡”族元素——具有铁磁性这一事实。我们可以理解为什么每个离子中的 d 电子倾向于组合成一个总自旋——在 Fe 中大到 $\frac{5}{2}$；却无法理解为什么相邻离子上的磁矩之间的相互作用会有利于平行取向。问题还因如下事实（§4.1）而更加复杂：这些所谓的 d 电子并非严格局域在特定的离子上，而是处于从原子到原子相互交叠、形成窄能带的状态。这个窄带又与普通的 s 带杂化，在 s 带中电子可以十分自由地传导。Heisenberg 模型哈密顿量 (10.29) 尽管作为金属铁磁性的唯象理论颇为成功，却需要从第一性原理出发加以认真的重新考察。这个微妙问题——固体理论中最重要的基本问题之一——的若干方面将在 §10.5–§10.6 中讨论。

\section{能带铁磁性}

从 Heisenberg 模型的局域自旋转向相反的极端，让我们把所有磁性电子都放进一个由 Bloch 态构成的能带。令 $n_{\mathbf{k}+}$ 为具有 + 自旋的态 $|\mathbf{k}\rangle$ 的占据数。在 Stoner 的\emph{集体电子模型}（collective electron model）中，我们在电子哈密顿量上添加如下一项：
\begin{equation*}
\mathcal{H}_{\mathrm{int.}}=\frac{U}{N}\sum_{\mathbf{k},\mathbf{k}'}n_{\mathbf{k}+}n_{\mathbf{k}'-}.
\tag{10.38}
\end{equation*}
于是，每一对自旋相反的电子都贡献一个大小为 $U/N$ 的正的“交换”能，它或许来自这两个电子偶尔同时栖身于同一原子的 d 壳层，但不一定恰好由像 (10.37) 那样的积分严格表示。自旋相同的电子的贡献则无须显式计入，因为它们可以归并到能量零点的定义之中。

虽然这一项可以解释为给出一个像 (10.22) 那样的内场，但我们现在处理的是 费米–狄拉克 分布，可以沿用 §10.2 的论证。令 $n_{\sigma}$ 为每原子中自旋为 $\sigma$ 的电子总数。代替 (10.9)，一个 + 自旋电子的能量为
\begin{equation*}
\mathcal{E}_{\mathbf{k}+}=\mathcal{E}(\mathbf{k})-\mu_{0}H+Un_{-},
\tag{10.39}
\end{equation*}
对 $-$ 自旋的电子有类似的式子。

两个电子集合的化学势仍然相同，因此我们必须满足如下方程：
\begin{equation*}
n\langle\mu\rangle=\mu_{0}\int_{0}^{\infty}\frac{1}{2}\bigl\{f^{0}(\mathcal{E}_{\mathbf{k}+})-f^{0}(\mathcal{E}_{\mathbf{k}-})\bigr\}\mathcal{N}(\mathcal{E})\,d\mathcal{E},
\tag{10.40}
\end{equation*}
以及
\begin{equation*}
n=\int\frac{1}{2}\bigl\{f^{0}(\mathcal{E}_{\mathbf{k}+})+f^{0}(\mathcal{E}_{\mathbf{k}-})\bigr\}\mathcal{N}(\mathcal{E})\,d\mathcal{E}.
\tag{10.41}
\end{equation*}
在 $H\to 0$、$T\to 0$ 的极限下，这些方程有一个与 (10.11) 类似的解。磁化率取如下形式：
\begin{equation*}
\chi=\frac{\mu_{0}^{2}\mathcal{N}(\mathcal{E}_{F})}{1-\frac{1}{2}U\mathcal{N}(\mathcal{E}_{F})}
\tag{10.42}
\end{equation*}
就好像我们只是把 Pauli 磁化率 (10.12) 增大了一个因子 $1/\{1-\frac{1}{2}U\mathcal{N}(\mathcal{E}_{F})\}$。交换相互作用 (10.38) 由于偏向平行自旋，使系统更容易被极化。

\begin{figure}[H]
\centering
\includegraphics[width=0.72\textwidth]{fig_177.pdf}
\caption*{图 177\quad (a) 能带铁磁性；(b) 铁磁金属的电子比热。}
\end{figure}

但是，当交换场强到足以使
\begin{equation*}
\frac{1}{2}U\mathcal{N}(\mathcal{E}_{F})>1,
\tag{10.43}
\end{equation*}
时，这个形式解显然是不稳定的。这标志着向铁磁性的转变。为使 (10.39)–(10.41) 相互一致，我们必须让 $n_{+}$（比如说）远大于 $n_{-}$，从而使系统具有一个永久磁矩。随后便可数值求解这些方程，给出 $\langle\mu\rangle$ 作为温度的函数。其结果在定性上与 (10.26) 的结果并无明显不同。系统表现出一个居里温度，在居里温度以下磁矩饱和到近乎恒定的值；在居里点以上，理论也预言了某种类似于居里–外斯定律 (10.24) 的行为。

在这个模型中，不难显式计算出与这一转变相联系的比热。能量（参见 §4.7）可以作为 $T$ 的函数算出，再对它求导便给出比热。结果表明，$C_{\mathrm{el.}}$ 在居里
点处出现不连续，紧挨其下方是一个尖锐的峰。这又是向铁磁性转变时观测到的一个特征，尽管理论并不能在所有细节上都重现该峰的确切形状。

铁磁性的巡游电子（itinerant electron）图像对铁磁金属绝大多数观测性质的解释，与局域自旋模型同样成功。它也与普遍的能带理论（§3.10）、输运理论以及费米面研究相一致；例如，在其中可以观察到两套分立的能带系统，分别容纳“向上”自旋与“向下”自旋的电子。另一方面，“交换”能 $U$ 的确切本性一点也不清楚，而且有独立的证据表明磁矩存在某种程度的表观局域化（§10.6、§10.11）。当前最好的看法是：对具有 d 带的过渡金属倾向于能带模型，而稀土金属中部分填充的 f 态里的自旋，则似乎总是表现得如同按 Heisenberg 方式局域化。

如 §8.3 所述，有些物质是铁电体（ferro-electric）：它们在没有电场的情况下也能携带永久电矩。这一现象在宏观上与铁磁性相似，但其来源是一个复杂得多、也极不寻常的微观效应。它出自某些晶体结构的一种能力：使自己发生轻微畸变，从而让晶格的每一个晶胞都带有一个偶极矩。

\section{磁性杂质}

过渡元素或稀土元素的原子溶入普通金属后，往往保留局域磁矩。这是一系列有趣的实验现象和理论概念的源泉。用相移语言（§5.5）所作的 Friedel 解释，与通过 Anderson 哈密顿量（Anderson Hamiltonian）的传统描述互为补充；借助后者，可以用少数几个唯象参量对这类现象的许多特征作出半定量估计。

首先我们假设：自旋相反的电子之间的“交换”或“关联”排斥，只有当二者处在同一个能量为 $\mathcal{E}_d$ 的原子 $d$ 能级上时才起作用。利用式 (10.38) 中已经用过的简化 Hartree--Fock 近似，我们写出

\begin{equation*}
\mathcal{H}_{dd} = U n_{d+} n_{d-} \tag{10.44}
\end{equation*}
其中 $n_{d\sigma}$ 表示自旋为 $\sigma$ 的 $d$ 电子的平均数目。倘若过渡元素溶入普通金属后这些原子能级仍保持完全锐利，它们无疑会被极化——比方说 $n_{d+}$ 取尽可能大的值，而 $n_{d-}$ 取尽可能小的值——这正是 §10.5 所讨论的那类情形的一个夸张版本。在这种组态下，各能级的有效能量被劈开，即（图 178）

\begin{equation*}
\mathcal{E}_{d\sigma} = \mathcal{E}_d + U n_{d\sigma}. \tag{10.45}
\end{equation*}
\begin{figure}[H]
\centering
\includegraphics[width=0.72\textwidth]{fig_178.pdf}
\caption*{图 178\quad 原子 $d$ 能级(a)被 $d$–$d$ 相互作用极化并劈裂(b)，随后因与 $s$ 电子杂化而展宽为共振峰(c)。}
\end{figure}

现在我们在这些局域 $d$ 态与溶剂金属传导电子（“s 电子”）的 Bloch 态之间引入一个强度为 $V$ 的相互作用。这个相互作用不一定对自旋敏感；原则上，它与过渡金属能带结构的模型哈密顿量表述（§3.10）中的 $s$–$d$ 杂化矩阵元基本相同，而那个矩阵元本可以由诸如式 (3.27) 那样的 L.C.A.O.“重叠积分”导出。因此我们假设这个相互作用只联系同种自旋的 $s$ 态和 $d$ 态，并且忽略 $s$ 电子自身之间可能存在的微弱交换相互作用。

然而，若把这个相互作用看作作用在每个 $d$ 能级上的微扰，它将引起与 $s$ 带之间的往返跃迁，其速率为

\begin{equation*}
W_\sigma/\hbar = (\pi/\hbar)\, V^2 \mathcal{N}_s(\mathcal{E}_{d\sigma}) \tag{10.46}
\end{equation*}

在这个“黄金定则”（Golden Rule）公式里，溶剂金属的性质仅通过其 $s$ 带中处于所考虑 $d$ 态能量 (10.45) 处的态密度体现出来。但这段被缩短的寿命可以解释为 $d$ 能级在能量上展宽了 $W_\sigma$。在 $\mathcal{E}_{d\sigma}$ 处的锐利能级不复存在，我们得到的是典型 Lorentz 线型的谱密度函数

\begin{equation*}
\rho_{d\sigma}(\mathcal{E}) = \frac{2l+1}{\pi}\, \frac{W_\sigma}{(\mathcal{E} - \mathcal{E}_{d\sigma})^2 + W_\sigma^2}. \tag{10.47}
\end{equation*}

原则上，这条谱线的中心也会稍有移动，但这不是一个重要的物理效应。

这个公式其实正是我们把 Friedel 求和 (5.40) 作为能量的函数、在诸如 (3.81) 那样的共振附近求导时应得的结果。因此，把宽度 $W_\sigma$ 解释为 (10.46) 中 $s$–$d$ 矩阵元 $V$ 的做法，并不是物理本质的要素，尽管在这些相互作用不太强时这些方程是有效的。

但是现在，一旦杂质上的 $d$ 能级被展宽，我们就不能再肯定它们仍会保持磁性极化。此时的情形其实与 (10.39)–(10.42) 完全一样，只需用 (10.47) 代替 $\frac{1}{2}\mathcal{N}(\mathcal{E})$。于是我们得到杂质离子上存在永久矩的一个条件：如同 (10.43)，若

\begin{equation*}
U \rho_{d+}(\mathcal{E}_F) > 1 \tag{10.48}
\end{equation*}
则极化态在能量上就是稳定的，不会有电子从 $n_{d+}$ 分布转移到 $n_{d-}$ 分布（参见 §5.2）。

过渡元素和稀土元素稀薄合金的多种可观测性质，都可以由这个理论得出，其中的参量 $U$ 和 $V$ 取合理而自洽的数值。例如，从 V 到 Co 的铁族元素溶解在 Cu、Ag、Au 中时是磁性的，但作为杂质出现在 Al 中时却不显示磁矩。Al 中传导电子浓度较高，使 (10.46) 中的 $\mathcal{N}_s$ 过大，从而使 $d$ 能级展宽，压低 (10.47) 中 $\rho_{d\sigma}(\mathcal{E})$ 的峰值，以致磁化条件 (10.48) 不再能得到满足。

基体金属的传导电子对杂质的存在并非无动于衷，只是这一效应依赖于自旋。设磁性离子完全极化，总自旋为 $S$。若所有“向上”自旋的 $d$ 态都已填满，同种自旋的 $s$ 电子便无法进入。因此这时只有“向下”自旋的 $s$ 电子受益。按普通微扰论估计，经由矩阵元 $V$ 虚跃迁到比费米能级高出能量 $U$ 的 $d$ 能级上，会使这些电子中每一个的能量改变

\begin{equation*}
J/N \approx -V^2/U. \tag{10.49}
\end{equation*}
于是系统的行为就好像存在如下形式的依赖于自旋的相互作用能

\begin{equation*}
\mathcal{H}_{sd} = -(J/N)\, \mathbf{s} \cdot \boldsymbol{\mathsf{S}}_i\, \delta(\mathbf{r} - \mathbf{R}_i) \tag{10.50}
\end{equation*}
杂质在 $\mathbf{R}_i$ 处的局域自旋矩 $\boldsymbol{\mathsf{S}}_i$ 与传导电子在 $\mathbf{r}$ 处的自旋 $\mathbf{s}$ 联系起来。注意，无论 $V$ 的符号如何，这一相互作用的符号都倾向于使两个自旋反平行（见 §10.7）。

这个 $s$–$d$ 相互作用 (10.50) 是作用在自由电子气上的一个高度局域化的微扰。为了领会它的效应，让我们计算系统的磁“响应”作为波数的函数，做法与我们在 §5.1 中计算静电响应函数基本相同。眼前的问题其实要简单得多，因为我们忽略了电子气内部的自旋–自旋相互作用，可以把整个效应都归因于“外”微扰。“向上”自旋和“向下”自旋的 $s$ 电子各带磁矩 $g\beta\mathbf{s}$，可以分开处理；按照导出 (5.16) 的论证，我们得到一个广义磁化率（参见 §8.1）：

\begin{equation*}
\chi_0(q) = 2(g\beta s)^2 \sum_{\mathbf{k}} \frac{f^0(\mathbf{k}) - f^0(\mathbf{k}+\mathbf{q})}{\mathcal{E}(\mathbf{k}+\mathbf{q}) - \mathcal{E}(\mathbf{k})}, \tag{10.51}
\end{equation*}
其中求和遍及两种自旋的电子态。

在自由电子模型中，这个和可以像 (5.35) 那样求出，结果是一个简单的解析函数，在 $q = 2k_F$ 处有二阶导数奇性。再作 Fourier 变换（参见式 (5.17)），便给出局域矩 (10.50) 邻域内磁化强度的空间分布。该奇性产生一个波数为 $2k_F$ 的振荡项，在远距离处按 $1/r^3$ 衰减（图 179）。这就是 R.K.K.Y.（Ruderman--Kittel--Kasuya--Yosida）效应；距杂质远近不同的各个离子，其 Knight 位移（§10.2）会有微小变化，可以用它来探测这一效应。

用 §5.5 的语言，我们也可以把磁性极化的杂质描述成一个区域：“向上”自旋电子因 d 相移中出现共振 (3.81) 而被吸引进去。而“向下”自旋电子的相应共振出现在高得多的能量处，因此两种分布在费米能级处的相移相差很大。由 (5.41)，这就给出磁极化的振荡，与静电杂质周围的电荷 Friedel 振荡完全类似。这种描述与 Anderson 模型等价，但不能像 (10.49) 那样直接给出用 $V$ 和 $U$ 表示的 $J$ 的数值。

或许更重要的是，这一现象为两个相邻磁性离子上磁矩之间的耦合提供了一种机制。“向上”自旋杂质紧邻处“向下”自旋 $s$ 电子的过剩，会在该邻域内的任何其他杂质上诱导出相应的“向上”自旋。

\begin{figure}[H]
\centering
\includegraphics[width=0.72\textwidth]{fig_179.pdf}
\caption*{图 179\quad $s$ 电子围绕杂质 A 的自旋极化 Friedel 振荡，导致与近邻杂质 B 的铁磁性相互作用，但在更远的距离（如 C 处）则可能有利于反铁磁性。}
\end{figure}

这样，我们就有了铁磁性的 Heisenberg 交换模型（§10.4）所需的全部机制：过渡金属离子 $d$ 壳层中局域化的磁矩，由于其极化方向倾向于与浸泡着它们的传导电子气反平行，而被促成为彼此平行。

这类相互作用在稀薄合金中引起复杂的磁学和热学效应，其中磁性离子是无规分布在溶剂晶格上的。在极端情形下，把非磁性溶剂去掉，留下纯过渡金属，这就是诸如 Fe 等金属中铁磁性的 Vonsovski--Zener 模型（Vonsovski--Zener model）。注意，这是一个比 Heisenberg--Dirac 哈密顿量 (10.29) 所提示的“动力学”得多的模型。$d$ 壳层磁矩并不是由整数个完全局域的“d 电子”组成，而是代表 $s$ 带与 $d$ 带已充分杂化的 Bloch 态中的平均贡献。因此，这个模型与 §10.5 中讨论的较简单的能带铁磁性并不矛盾，只是参量 $U$ 需要重新解释。

依赖于自旋的 $s$–$d$ 相互作用 (10.50) 还导致一个与磁性杂质对电阻率的贡献相联系的奇特效应。设我们像 (6.90) 中那样，试图计算这样一个物体把一个处于态 $|\mathbf{k},+\rangle$ 的“向上”自旋 $s$ 电子散射到同自旋态 $|\mathbf{k}',+\rangle$ 的微分截面。

在第一级 Born 近似下，散射振幅恰好就是矩阵元

\begin{align*}
t^{(1)} &= \langle \mathbf{k}, + | \mathcal{H}_{sd} | \mathbf{k}', + \rangle \notag \\
&= -(J/N)\, \mathsf{S}_i^{(z)}. \tag{10.52}
\end{align*}
这一项平方后不为零，而且是价电荷散射等（§7.4）之外的附加项，本来通常就足以满足我们对此问题的好奇心了。

\begin{figure}[H]
\centering
\includegraphics[width=0.72\textwidth]{fig_180.pdf}
\caption*{图 180\quad 二级 $s$–$d$ 散射中的直接过程(a)与交换过程(b)。注意自旋翻转跃迁中 $s$ 自旋与 $d$ 自旋的相对取向。}
\end{figure}

但是，本着一种迂腐的精神，让我们考察微扰级数中更高阶的项。初等量子力学教科书教导我们，第二级 Born 近似取如下形式

\begin{equation*}
t^{(2)} = \sum_{\mathbf{k}'',\sigma} \frac{\langle \mathbf{k}, + | \mathcal{H}_{sd} | \mathbf{k}'', \sigma \rangle \langle \mathbf{k}'', \sigma | \mathcal{H}_{sd} | \mathbf{k}', + \rangle}{\mathcal{E}(\mathbf{k}) - \mathcal{E}(\mathbf{k}'')}, \tag{10.53}
\end{equation*}
其中“中间”态 $|\mathbf{k}'',\sigma\rangle$ 通过相互作用哈密顿量 $\mathcal{H}_{sd}$ 与初态和末态两者相联系。然而，这个表达式并不完全正确。在完整的时间依赖微扰论中，我们必须考虑两类“图”（图 180）。在直接
过程（a）中，入射电子先进入态 $|\mathbf{k}'',\sigma\rangle$（因此该态必须是空的），然后才被发射到末态；而在交换（exchange）过程（b）中，我们假设一个已处于 $|\mathbf{k}'',\sigma\rangle$ 的电子，在“初始”电子被接纳进中间态之前，就先以 $|\mathbf{k}',+\rangle$ 离去。由这些规则，我们应当把 (10.53) 改写成

\begin{align*}
t^{(2)} = \sum_{\mathbf{k}'',\sigma} \frac{1}{\mathcal{E}(\mathbf{k}) - \mathcal{E}(\mathbf{k}'')} \big\{ &(1 - f_{\mathbf{k}''}) \langle \mathbf{k}, + | \mathcal{H}_{sd} | \mathbf{k}'', \sigma \rangle \langle \mathbf{k}'', \sigma | \mathcal{H}_{sd} | \mathbf{k}', + \rangle \\
&+ f_{\mathbf{k}''} \langle \mathbf{k}'', \sigma | \mathcal{H}_{sd} | \mathbf{k}', + \rangle \langle \mathbf{k}, + | \mathcal{H}_{sd} | \mathbf{k}'', \sigma \rangle \big\}, \tag{10.54}
\end{align*}
其中 $f_{\mathbf{k}''}$ 是中间态的占据数。

对于普通的势散射，(10.54) 中两个矩阵元的乘积与其次序无关，于是 $f_{\mathbf{k}''}$ 相消，我们便回到 (10.53)，它对散射截面的贡献只是对一级项 (10.52) 的一个小的修正。但对自旋散射我们必须记住，$\mathcal{H}_{sd}$ 中含有诸如 $\mathsf{S}_i^{(x)} s_x$ 和 $\mathsf{S}_i^{(y)} s_y$ 这样的项，它们以消耗杂质的极化为代价去“翻转”传导电子的自旋（见 (10.104)）。算符 $\mathsf{S}_i^{(x)}$ 与 $\mathsf{S}_i^{(y)}$ 不对易，因此 (10.54) 中各项的先后次序变得重要。从杂质自旋矩的角度看，两个跃迁中哪一个先发生确实至关重要。例如显而易见，倘若杂质已经具有最大的“向上”自旋，直接过程就是被禁戒的。

结果是，在 (10.54) 的代数约化中最终留下一个如下形式的项

\begin{equation*}
t_{\rm K} \approx 2(J/N)^2 \mathsf{S}_i^{(z)} \sum_{\mathbf{k}''} \frac{f_{\mathbf{k}''}}{\mathcal{E}(\mathbf{k}) - \mathcal{E}(\mathbf{k}'')}, \tag{10.55}
\end{equation*}
其中因子 $f_{\mathbf{k}''}$ 并未被来自其他图的贡献所抵消。由于 $f_{\mathbf{k}''}$ 在费米能量处陡然变化，当 $\mathcal{E}(\mathbf{k})$ 趋近 $\mathcal{E}_F$ 时，这个和会变得很大。例如，假设传导带中的态密度一直低到某个能量 $(\mathcal{E}_F - D)$ 都具有 $\mathcal{N}(\mathcal{E}_F)$ 的量级，我们可以把 (10.55) 近似地算成

\begin{equation*}
t_{\rm K} \approx 2(J/N)^2 \mathsf{S}_i^{(z)} \mathcal{N}(\mathcal{E}_F) \ln \frac{D}{|\mathcal{E}_F - \mathcal{E}(\mathbf{k})|}. \tag{10.56}
\end{equation*}
普通带电杂质在金属中产生的\emph{剩余电阻}（residual resistance）与温度无关（§§7.2--7.4）。但是，当我们把磁性杂质引起散射的总概率在费米面上厚度为 $kT$ 的“热层”内积分时，
(10.56) 中的奇性并未因此被抹平。电阻率应当随温度作对数式变化，大致为

\begin{equation*}
\rho(T) \sim \rho_0 \{1 - 2(J/N)\,\mathcal{N}(\mathcal{E}_F) \ln D/kT\}. \tag{10.57}
\end{equation*}

由于 $J$ 本质上是负的，这个公式预言剩余电阻率随温度下降而迅速增大。再把它与基体金属低温下按 $T^5$ 变化的“理想”电阻（§7.5）结合起来，就解释了长期在含过渡元素的稀薄合金中观测到的神秘的电阻极小（resistance minimum）。而 (10.56) 的强能量依赖进入诸如 (7.104) 那样的公式，则解释了在这类材料中同样观测到的巨大温差电功率（giant thermo-electric power）。

但完整的 Kondo 效应（Kondo effect）理论还必须考虑近邻杂质的相互极化、$t$ 的微扰展开的高阶修正、具有 $s$–$d$ 相互作用的系统之基态的性质，以及许多其他微妙之处。这个领域如今已成为固体物理中各种高深量子力学方法的主要演练场之一。

\begin{figure}[H]
\centering
\includegraphics[width=0.72\textwidth]{fig_181.pdf}
\caption*{图 181\quad (a)铁磁有序。(b)反铁磁有序。}
\end{figure}

\section{反铁磁性}

Heisenberg 模型中铁磁性的一般条件是 (10.29) 中的交换积分 $J_{ll'}$ 为正。此时内场与自旋的排列方向平行，整个系统倾向于取图 181(a) 的组态。

但假如——正如对 $J_{ll'}$ 的计算常常表明的那样——该积分的符号为负，那么任一格点上的内场 (10.32) 就会与近邻自旋磁化的平均方向相反，图 181(b) 的排列成为有利。我们把这样的系统称为反铁磁性（antiferromagnetic）的。

为了看出会发生什么，让我们分别考虑自旋取向相反的两套子晶格。设
$\boldsymbol{\mu}_+$ 是 $\oplus$ 子晶格上原子的平均磁化。交换相互作用将在 $\ominus$ 子晶格上产生一个内场（不妨写成）

\begin{align*}
\mathbf{H}^- &= \frac{1}{\beta^2}(\Sigma J)\boldsymbol{\mu}_+ \notag \\
&= \lambda N \boldsymbol{\mu}_+, \tag{10.58}
\end{align*}
这个求和假定只对最近邻进行，因为 $J_{ll'}$ 被认为作用程极短。要点在于，$\lambda$ 现在将是\emph{负}的。

同样由于这一相互作用，每个 $\oplus$ 格点上也存在内场

\begin{equation*}
\mathbf{H}^+ = \lambda N \boldsymbol{\mu}_-. \tag{10.59}
\end{equation*}

现在假设我们的磁体是经典的，即 (10.17) 对磁场 $\mathbf{H}$ 中磁矩 $\boldsymbol{\mu}$ 的\emph{平均}值成立。把 (10.17) 写成

\begin{equation*}
\langle \boldsymbol{\mu} \rangle = \mathcal{L}\{(\boldsymbol{\mu}\cdot\mathbf{H})/kT\}, \tag{10.60}
\end{equation*}
其中用符号 $\mathcal{L}$ 表示由 (10.17) 右端定义的 Langevin 函数（Langevin function）。于是我们有两个方程要求解：

\begin{align*}
\boldsymbol{\mu}_- &= \mathcal{L}\{\boldsymbol{\mu}\cdot(\mathbf{H}+\lambda N\boldsymbol{\mu}_+)/kT\}, \notag \\
\boldsymbol{\mu}_+ &= \mathcal{L}\{\boldsymbol{\mu}\cdot(\mathbf{H}+\lambda N\boldsymbol{\mu}_-)/kT\}. \tag{10.61}
\end{align*}

在高温下，当近似 (10.18) 有效时，这两个方程可以相加而给出 (10.23)。结果与先前的 (10.24) 几乎完全一样：

\begin{equation*}
\chi = \frac{N\langle\mu^2\rangle}{3k(T+\theta)}; \tag{10.62}
\end{equation*}
$\lambda$ 为负时，方便的做法是按 $-\lambda N\langle\mu^2\rangle/3kT$ 来定义 $\theta$，这与 (10.25) 相一致。磁矩倾向于彼此反平行地排列，而不是被磁场拉向同一方向，顺磁磁化率因此减小。

但在低温下，即使不存在磁场也有一个解。若我们写出

\begin{equation*}
\boldsymbol{\mu}_- = -\boldsymbol{\mu}_+ \tag{10.63}
\end{equation*}
（在图 181(b) 的有序反铁磁排列中，直觉上正是如此），则 $\lambda$ 的负号被抵消，我们得到方程

\begin{equation*}
\boldsymbol{\mu}_+ = \mathcal{L}\{(\boldsymbol{\mu}\cdot|\lambda|N\boldsymbol{\mu}_+)/kT\} \tag{10.64}
\end{equation*}
待解。这个方程与 §10.3 中导致铁磁性的方程完全相同——例如 (10.26)，其中含有 Langevin 函数的“量子”等价物。

这样，从单个自旋的角度看，铁磁性与反铁磁性似乎很相似：系统在高温下是顺磁性的，但在某一温度以下变成强极化；在反铁磁情形中，这个温度称为 Néel 温度（Néel temperature），

\begin{equation*}
T_{\rm N} = \theta. \tag{10.65}
\end{equation*}

然而在宏观上，两者的行为大不相同。反铁磁排列不具有任何净极化，因此系统不显示任何永久磁化。所能观测到的只是磁化率的变化；要计算它，需在一个小的磁场 $H$ 之下求解 (10.61)。

\begin{figure}[H]
\centering
\includegraphics[width=0.72\textwidth]{fig_182.pdf}
\caption*{图 182\quad 反铁磁体的磁化率：(a)$\chi_\parallel$ 趋于零；(b)$\chi_\perp$ 保持恒定；(c)$\chi = \frac{1}{3}\chi_\parallel + \frac{2}{3}\chi_\perp$。}
\end{figure}

这个计算原则上很简单，只是实际做起来略嫌繁琐。其结果有趣之处在于，它依赖于 $\mathbf{H}$ 相对于两套子晶格磁化方向的取向。当 $\mathbf{H}$ 与 $\boldsymbol{\mu}_+$ 平行时，磁化率 $\chi_\parallel$ 在 $T=0$ 处趋于零。这是因为内场强到足以阻止 $\ominus$ 子晶格中的任何自旋翻转。另一方面，$\chi_\perp$ 在 $T=0$ 的值与它在 $T_{\rm N}$ 处应有的值相同：与外场垂直的子晶格极化，对其磁化被稍微转向外场方向只有微弱的抗拒。

实际上，一个宏观样品将包含大量磁畴，各磁畴的极化轴彼此不同，于是对平均磁化率可以写出

\begin{equation*}
\chi = \tfrac{1}{3}\chi_\parallel + \tfrac{2}{3}\chi_\perp. \tag{10.66}
\end{equation*}
$\chi$ 在 $T_{\rm N}$ 处斜率的不连续性很容易观测到，并且与子晶格极化时出现的比热反常相联系。

值得指出的是，由于 $\chi_\perp > \chi_\parallel$，磁场本应能够把所有磁畴都拉成垂直组态。但这种情况并不发生——大概是由于自旋–自旋相互作用的各向异性，它倾向于偏爱某些特殊的晶体方向。

我们在这里考虑的模型，只是展现着令人眼花缭乱的种种现象的一大类系统中的一个简单例子。这个论证可以从不同方向加以细化。例如，晶体结构可能是复杂的，其中磁性离子并不形成简单立方网络，系统究竟会如何划分成子晶格也可能并非一目了然。又如，相互作用的作用程可能比最近邻间更长，或者对某些自旋对是反铁磁性的，而对另一些则是铁磁性的。这样做的一个后果（在我们的框架内很容易处理：只需在 (10.58) 中加上诸如 $\lambda' N\boldsymbol{\mu}_-$ 的项，在 (10.59) 中加上 $\lambda' N\boldsymbol{\mu}_+$ 的项）是改变 Néel 温度，使之不再与 (10.62) 中的参量 $\theta$ 相同。

\begin{figure}[H]
\centering
\includegraphics[width=0.72\textwidth]{fig_183.pdf}
\caption*{图 183\quad 亚铁磁性。}
\end{figure}

\begin{figure}[H]
\centering
\includegraphics[width=0.72\textwidth]{fig_184.pdf}
\caption*{图 184\quad 超交换。}
\end{figure}

另一种情形是晶体中存在两种不同类型的磁性离子，而且二者不具有相同的磁矩。于是我们可能遇到图 183 所描绘的情况：A、B 两类原子的自旋取向相反，但由于一套子晶格上的矩比另一套上的大，系统仍将显示永久矩。这称为亚铁磁性（ferrimagnetism），是铁氧体（ferrite）的特征；在宏观上，它看起来与铁磁性非常相像。

在许多反铁磁物质中，磁性离子彼此相距相当远，以致很难相信它们能由诸如 (10.37) 那样的交换积分联系起来。这并不意味着晶格非常稀疏，而是意味着在磁性离子之间的空隙里填有许多非磁性离子。例如，它们可以是氧离子，如 MnO 中那样。这时我们就遇到了称为超交换（super-exchange）的机制：Mn 离子中的 $d$ 电子与氧中的电子相耦合，而后者再通过交换作用与第二个 Mn 离子中的 $d$ 电子相互作用。

示意性地，情形可以如图 184 所示。若原子 A 具有 $+$ 自旋，就有空位让一个 $-$ 自旋的电子从氧的 $p$ 态进入。若原子 B 具有 $-$ 自旋，它就能接纳这对电子中来自氧的另一个。这样，氧上 $p$ 电子自旋之间的相互作用就转化为磁性离子 $d$ 壳层自旋之间的相互作用。这一效应的实际计算颇为微妙，但它至少在原理上解释了观测到的现象。

在稀磁合金中，§10.6 的间接交换（indirect exchange）机制也可以是反铁磁性的。如图 179 所示，当第二个杂质恰好处在与第一个磁性离子的 R.K.K.Y. 振荡中“向上”自旋峰相对应的距离上时，就会出现这种情形。但要理解诸如 Cr 这类纯金属中的反铁磁性，我们应当回到 §10.5 的巡游电子图像。

定量地看，不难理解反铁磁图样如何能被传导电子之间的交换作用所稳定。设图 181(b) 那种类型的磁有序已经建立。一个“向上”自旋的电子穿越晶体时，通过交换相互作用 $U$ 将感受到一个依赖于局域极化的有效场。这个场具有磁超晶格的周期性，因而会在通常布里渊区内部的一个“子区”边界上引起能量不连续（参见 §2.6）。能带结构中的这一特征会反映为“向上”自旋电子密度的一个超晶格。对于自旋相反的电子，微扰的符号相反，于是它们倾向于占据另一套超晶格。如果这些分布碰巧与我们最初设想的磁化图样相一致，我们就找到了电子的一种自洽的反铁磁有序，而无须强迫它们局域化。

为了给这类论证赋予数学实质，我们假设
我们把均匀磁化率 (10.42) 推广为一个依赖于 $q$ 的磁响应函数。沿着 §5.1 的思路，在导致磁化率 $\chi_0(\mathbf{q})$（即式 (10.51)）的那类计算中，纳入一个形式如 (10.38) 的交换相互作用是很容易的。所得结果实际上与式 (10.42) 明显相似，即

\begin{equation*}
\chi(\mathbf{q}) = \frac{\chi_0(\mathbf{q})}{1 - \tfrac{1}{2} U (g\beta s)^{-2} \chi_0(\mathbf{q})}.
\tag{10.67}
\end{equation*}

现在，如果这个表达式的分母为零，金属的假定非磁态便是不稳定的。磁有序的 Overhauser 条件是

\begin{equation*}
\tfrac{1}{2} U (g\beta s)^{-2} \chi_0(\mathbf{q}) > 1
\tag{10.68}
\end{equation*}

其中铁磁性的条件 (10.43) 正是 $q = 0$ 的特殊情形。

通常的自由电子间交换相互作用能否在普通金属中导致稳定的自旋密度波（spin density waves），这个问题尚未得到确定的解决。然而在过渡金属中，已知 $U$ 相当大，而待代入 $\chi_0(\mathbf{q})$ 的公式 (10.51) 的电子能态 $\mathcal{E}(\mathbf{k})$ 由能带结构主导。例如，由 §5.4 显然可知，一个横跨费米面相互平行的平坦区域的波矢 $\mathbf{q}$ 可能诱导出很大的 $\chi_0(\mathbf{q})$ 值，从而使 (10.68) 得以满足。在这种情形下，反铁磁有序可能异常复杂：自旋密度的变化等价于锥面上自旋矩的螺旋形排布，其重复距离与底层晶格无关。

\section{Ising 模型}

在讨论铁磁性和反铁磁性时，我们对 Heisenberg 哈密顿量 (10.29) 的态采用了非常粗糙的近似。实际上，我们假定除正在处理的那个自旋之外，每个自旋都取其平均值；我们忽略了相邻格点上自旋之间的涨落与关联。此外，我们还忽略了 $\mathsf{S}_l$ 是算符这一事实——它们不能换成各自的平均期望值。

自旋哈密顿量的本征值分布实际上非常复杂，试图对系统的热力学给出精确的讨论，是固体理论的主要课题之一。这一理论沿两个方向发展。我们可以假定系统接近其基态（铁磁的或反铁磁的），然后讨论低能激发。这就是自旋波（spin waves）理论，将在 §10.11 中概述。

另一类近似是略去自旋算符的非对角元，只考虑它们沿某个固定方向——通常是 $H$ 的方向——的分量。换言之，我们假定系统的能级由下式给出

\begin{equation*}
\mathcal{E} = -\sum_{ll'} J_{ll'} \sigma_l \sigma_{l'} - \beta H \sum_l \sigma_l,
\tag{10.69}
\end{equation*}

\begin{figure}[H]
\centering
\includegraphics[width=0.72\textwidth]{fig_185.pdf}
\caption*{图 185\quad Ising 模型的一个组态 (a) 可以代表：(b) 自旋的一种排布；(c) 二元合金中原子的一种排布；(d) "格气"的一个组态。}
\end{figure}

其中我们在每个格点 $l$ 上定义一个只能取 $+1$ 或 $-1$ 的“自旋”量子数 $\sigma_l$。在多数情况下我们还假定 $J_{ll'}$ 只在格子中最近邻之间起作用。

这就叫作 Ising 模型。它不可能是铁磁或反铁磁系统的精确描述，但由于它绕开了求自旋集合的精确本征函数这一极为困难的问题，使人们得以集中处理这类系统的态的统计分布这一组合问题。

Ising 模型确实更精确地描述了另一类系统——合金：两种元素 A 和 B 可以在一组固定的格点位置上相互替换。我们可以说 $\sigma_l = +1$ 对应于第 $l$ 个格点上的一个 A 原子，而 $\sigma_l = -1$ 规定该格点上是一个 B 原子。“交换积分” $J$ 这时应解释为

\begin{equation*}
J = \tfrac{1}{2}\bigl\{ \mathcal{V}_{AB} - \tfrac{1}{2}(\mathcal{V}_{AA} + \mathcal{V}_{BB}) \bigr\},
\tag{10.70}
\end{equation*}

其中 $\mathcal{V}_{AB}$ 是相邻格点上 A 原子与 B 原子之间的相互作用能，减去的则是一个 AA 对或一个 BB 对的平均能量。

这类合金系统可以作实验研究，而且发现它们的确表现出一个过渡温度，类似于居里温度或 Néel 温度。在转变点以上，合金是无序的；温度较低时它趋于有序，或者是通过分凝出近乎纯 A 与近乎纯 B 的材料（这对应于 $J > 0$ 的情形，即铁磁性），或者是通过形成 A、B 原子规则地分布在相互穿插的子晶格上的区域——这是反铁磁性的类似物。有序–无序转变（order–disorder transition）与铁磁或反铁磁转变的差别在于：前一情形中 A 类与 B 类原子的总数是固定的，而“向上”与“向下”的自旋却可以自由地相互转变。这一差别在理论上可以形式地加以顾及，因此许多结果对两类系统普遍成立。

Ising 模型可以应用的另一个问题是“格气”（lattice gas）或“格液体”（lattice liquid）问题：把一些能在短距离上相互作用的原子放到一组格点上，并留下空格点。这在形式上与合金问题等价。然而，它是否真正贴切地适用于人们有时让它去描述的液–气转变的热力学，却令人怀疑。

不过，Ising 模型的重要性并不在于描述个别的物理效应；它是一个数学上易于处理的模型，描写一个应当表现出合作现象（co-operative phenomena）和相变的系统。在这个意义上，它更多地属于统计力学的一般理论，而不属于固体理论。然而，有某些原理对晶体格子的一般理论也很重要，我们现在就试图阐明它们。

\section{组合方法}

Ising 模型的优点是，其热力学性质的计算可以归结为一个组合问题。例如，假设我们要求配分函数（partition function）

\begin{equation*}
Z = \sum e^{-\mathcal{E}/kT},
\tag{10.71}
\end{equation*}

其中求和遍及所有“组态”，即遍及格子的 $N$ 个格点上 $+$ 自旋与 $-$ 自旋的各种排布。这个和看起来非常复杂，但形式上它只需涉及两个变量，例如 $N_A$——“向上”自旋[或 A 原子]的数目——与 $N_{AB}$——相邻反平行自旋[或 AB 对]的数目。

为看清这一点，我们把能量 (10.69) 用各种类型的键的总数表出：

\begin{equation*}
\mathcal{E} = J(N_{AB} - N_{AA} - N_{BB}) - \beta H (N_A - N_B).
\tag{10.72}
\end{equation*}

但这些变量并非相互独立。显然有

\begin{equation*}
N_A + N_B = N.
\end{equation*}

同样显然的是，止于 A 格点的键至多只有 $pN_A$ 条，这里 $p$ 是格子的配位数（co-ordination number），即给定格点的最近邻数目。每条这样的键需要在 $N_{AA}$ 中被计及两次、在 $N_{AB}$ 中被计及一次。于是 $2N_{AA} + N_{AB} = pN_A$；类似地 $2N_{BB} + N_{AB} = pN_B$。代入这些条件，我们发现 (10.72) 变成

\begin{equation*}
\mathcal{E} = -\tfrac{1}{2} pNJ + 2N_{AB} J - (2N_A - N)\beta H.
\tag{10.73}
\end{equation*}

现在设 $g(N; N_A, N_{AB})$ 表示把 $N_A$ 个“向上”自旋放到格子上、且恰有 $N_{AB}$ 对反平行近邻的全部放法数。所有这些组态都具有相同的能量；我们可以把配分函数 (10.71) 写成

\begin{align*}
Z &= y^{\frac{1}{2}N} z^{-\frac{1}{2}pN} \sum_{N_A, N_{AB}} g(N; N_A, N_{AB})\, y^{N_A} z^{N_{AB}} \\
&= y^{\frac{1}{2}N} z^{-\frac{1}{2}pN} \Lambda_N(y, z),
\tag{10.74}
\end{align*}

其中变量 $y$ 和 $z$ 与磁场及相互作用 $J$ 相联系：

\begin{equation*}
y \equiv e^{2\beta H/kT}; \quad z \equiv e^{-2J/kT}.
\tag{10.75}
\end{equation*}

系统的全部有趣性质现在都集中在函数 $\Lambda_N$ 之中；由它可以导出适用于“铁磁”与“反铁磁”系统、或合金相与正规溶液（regular solutions）理论的公式。

但是，计算 (10.74) 中组合因子的问题一般是不可解的；已知的只是近似公式。对给定的 $N_A$ 值，算出组态的总数并不困难。由初等代数，这必是把 $N_A$ 个物体放到 $N$ 个位置上的排布方式总数

\begin{equation*}
\sum_{N_{AB}} g(N; N_A, N_{AB}) = \frac{N!}{N_A! (N - N_A)!}.
\tag{10.76}
\end{equation*}

困难在于把这个总数按 $N_{AB}$ 的各个取值分开。

一个非常粗糙的近似，是把对 $N_{AB}$ 的单独求和换成假定它取其平均值。也就是说，我们把 $N_A$ 个 A 原子和 $N_B$ 个 B 原子无规地分布在格子的格点上，并计算 $\tfrac{1}{2}pN$ 条键中任一条一端为 A 原子、另一端为 B 原子的概率。结果当然是

\begin{align*}
\langle N_{AB} \rangle &= 2\, \frac{N_A}{N}\, \frac{N_B}{N}\, \tfrac{1}{2} pN \\
&= \frac{p N_A N_B}{N}.
\tag{10.77}
\end{align*}

把 (10.76) 和 (10.77) 代入 (10.74)，得

\begin{equation*}
\Lambda_N \approx \sum_{N_A} \frac{N!}{N_A! N_B!}\, y^{N_A} z^{p N_A N_B / N}.
\tag{10.78}
\end{equation*}

由此可以计算自由能：从对数中取出最大的项，并用 Stirling 公式作渐近展开。于是

\begin{align*}
F &= -kT \ln \Lambda_N \\
&\approx kT \Bigl\{ -N \ln N + N_A \ln N_A + N_B \ln N_B - N_A \ln y - p\, \frac{N_A N_B}{N} \ln z \Bigr\}.
\tag{10.79}
\end{align*}

\begin{figure}[H]
\centering
\includegraphics[width=0.72\textwidth]{fig_186.pdf}
\caption*{图 186\quad 对 $N=5$、$N_A=3$ 的链作组态计数。此系统的 $\langle N_{AB}\rangle = 2.4$。}
\end{figure}

自由能取极小（也即求和中出现极大项）的条件，可通过对 $N_A$ 求导得到（当然要取 $N_B = N - N_A$）。利用 (10.75)，我们得到

\begin{equation*}
\ln N_A - \ln N_B - \frac{2\beta H}{kT} - \frac{(N_A - N_B)}{N}\, \frac{2pJ}{kT} = 0,
\end{equation*}

即

\begin{equation*}
\frac{N_A}{N_B} = \exp \Bigl\{ \frac{2\beta H}{kT} + \Bigl( \frac{N_A - N_B}{N} \Bigr) \frac{2pJ}{kT} \Bigr\},
\tag{10.80}
\end{equation*}

或

\begin{equation*}
\Bigl( \frac{N_A - N_B}{N} \Bigr) = \tanh \Bigl\{ \frac{\beta H}{kT} + \Bigl( \frac{N_A - N_B}{N} \Bigr) \frac{pJ}{kT} \Bigr\}.
\tag{10.81}
\end{equation*}

如果把 $(N_A - N_B)/N$ 认作 $\langle \mu \rangle$，即系统的平均净极化，并像 (10.33) 那样把 $pJ$ 表为一个内场，那么我们就回到了铁磁转变的居里–外斯条件 (10.26)。在合金理论中，这被称为 Bragg--Williams 近似（Bragg--Williams approximation）。

\begin{figure}[H]
\centering
\includegraphics[width=0.72\textwidth]{fig_187.pdf}
\caption*{图 187\quad 准化学近似下的比热，格子的配位数为 p = 12。}
\end{figure}

上面这个初等内场公式的推导暴露了它的局限性；它对相邻自旋之间的关联只字未提。例如在铁磁体中，由于自旋对有平行排列的倾向，我们可以预期随机平均值 (10.77) 高估了实际平衡组态中 AB 对的数目。

改进这一点的一个有趣途径是 Guggenheim 的准化学近似（quasi-chemical approximation）。假设我们在无规混合的基础上算出与 (10.77) 类似的其他平均值，就得到

\begin{equation*}
\frac{\langle N_{AB} \rangle^2}{\langle N_{AA} \rangle \langle N_{BB} \rangle} = 4.
\tag{10.82}
\end{equation*}

但如果我们处理的是分子之间的化学平衡，如同化学反应

\begin{equation*}
AA + BB \rightleftharpoons 2AB,
\tag{10.83}
\end{equation*}

所表达的那样，那么我们知道，不同组元的平衡浓度应由下式给出

\begin{equation*}
\frac{\langle N_{AB} \rangle^2}{\langle N_{AA} \rangle \langle N_{BB} \rangle} = 4\, \frac{\bigl( e^{-\mathcal{E}_{AB}/kT} \bigr)^2}{e^{-\mathcal{E}_{AA}/kT}\, e^{-\mathcal{E}_{BB}/kT}},
\tag{10.84}
\end{equation*}

其中 $\mathcal{E}_{AB}$ 等等是 AB 分子等等的能量。用 (10.75) 的记号，这提示 (10.84) 的右端平均而言应等于 $4z^2$。对 (10.73) 稍作变换，便可算出 $\langle N_{AB} \rangle$，并用它代替 (10.77) 放入 (10.78)。

这个计算我们就不再往下做了，只想指出：在由下式定义的转变温度以下，系统表现出自发极化（铁磁性或反铁磁性）

\begin{equation*}
\frac{kT_c}{pJ} = \frac{2/p}{\ln(1 - 2/p)},
\tag{10.85}
\end{equation*}
（译注：按准化学近似的标准结果，右端应为 $-2/p\big/\ln\frac{p}{p-2}$，原书疑漏负号。）

在配位数 $p$ 很大的极限下，它应趋于 (10.28)（在我们现在的记号中 $k\theta = pJ$）。我们还可以近似地计算转变温度附近的比热；它表现出一个尖锐的奇点，在高温一侧有一间断，如图 187 所示。

碰巧，准化学近似等价于另一个看上去颇不相同的方法，它出自 Bethe 和 Peierls。这是把内场论证加以推广的一种尝试。取一个自旋集团——例如中心一个位于 $l = 0$ 的自旋，以及它的 $p$ 个近邻。从这个集团的能量中取 (10.69) 的下列各项：

\begin{equation*}
\mathcal{E}_p = -\sum_{j=1}^{p} J \sigma_0 \sigma_j - \beta H \sigma_0 - \beta H_I \sum_{j=1}^{p} \sigma_j.
\tag{10.86}
\end{equation*}

换言之，我们假定中心自旋感受到实际的外场 $H$，并且显式地与近邻的自旋相连；而对集团的外围格点，我们则假定可以把它们与环境的相互作用当作一个内场 $H_I$ 的效应来处理。

\begin{figure}[H]
\centering
\includegraphics[width=0.72\textwidth]{fig_188.pdf}
\caption*{图 188\quad Bethe–Peierls 集团。}
\end{figure}

很容易为这个集团写出完整的配分函数，因为它只含 $p$ 项。于是可以算出集团中心原子的平均取向。但中心格点并不是格子中什么特权位置；$\langle \sigma_0 \rangle$ 必须与 $\langle \sigma_j \rangle$ 相同，后者是集团边缘上一个自旋的平均矩，它同样可以写成 $H_I$ 的函数。这就给出了内场 $H_I$ 的一个自洽条件：即使 $H = 0$，在低于由 (10.75) 定义的转变温度 $T_c$ 的温度下，$H_I$ 也可以不为零。

Bethe 方法作为一种物理上直观的程序是有吸引力的，甚至可以让它对更现实的模型给出结果：引入真正的自旋算符来代替 (10.86) 的“Ising 自旋变量”，并对这个集团哈密顿量的能级计算配分函数。或者，我们也可以在第一近邻的简单集团之外再引入更多的壳层，把问题解到更高的近似。

这些更高阶的近似，实际上是想更准确地确定 $\Lambda_N(y, z)$ 幂级数展开的系数。可以证明，每一个系数——无论是低温下有效的按 $z$ 的幂的展开，还是高温下有效的按

\begin{equation*}
w = \tanh \frac{J}{kT},
\tag{10.87}
\end{equation*}

的幂的展开——都依赖于数出在格子的“键”上能够描出的、具有给定边数的多边形，这颇像非理想气体理论中的 Mayer 集团和展开，或微扰论中的“图形”。这种计数原则上足够容易，实际做起来却极其繁难。

此外，这样的展开还不能完全使人相信：系统在某个确定温度的热力学性质中真的存在奇点；而 Bragg--Williams 方法和 Bethe 方法的封闭解析公式也填补不了这个空缺。这些方法\emph{假定}存在一个内场，或者假定系统有一个平均说来在每个格点上都相同的净极化。换言之，它们假定系统具有长程序（long-range order）——也就是说，只要知道一个格点上自旋的方向，我们就能猜出任何其他格点上它的方向，不管有多远。\footnote{在铁磁情形，要知道整个格子极化的净方向，只需知道一个自旋。在反铁磁情形，远处格点上自旋的方向当然还取决于它是否与原点的格点属于同一个子晶格。}

\begin{figure}[H]
\centering
\includegraphics[width=0.72\textwidth]{fig_189.pdf}
\caption*{图 189\quad (a) 长程序蕴含短程序。(b) 短程序并不蕴含长程序。}
\end{figure}

我们前面已经指出，Bragg--Williams 方法甚至给不出任何短程序（short-range order），除非它是被假定的长程序所蕴含的：在居里温度以上，分布被假定得完全无规，每个自旋都与近邻的详细行为无关。Bethe 方法在 $T_c$ 以上会给出短程序，因此更为现实。两种方法都能相当好地描述长程有序态在充分建立之后的一般性质，但都不能用来证明这个态的确存在。

\section{Ising 问题的严格解}

事实证明，Ising 问题可以在两种特殊情形下严格求解。其中一种情形是平庸的，证明也很简单；另一种情形表述起来并不困难，但配分函数最终公式的证明却太过艰难，无法在这里给出。

平庸的情形是线性链；我们给出它的解，因为它展示了一种普遍的技巧。设考虑由 $N$ 个“层”组成的格子。在一维情形，这样一个“层”就是单个格点；在二维情形，它会是完整的一行原子，依此类推。

现在设 $\nu_j$ 是标志第 $j$ 层“状态”的记号；例如在二维情形，$\nu_j$ 可以是 $(++-+\ldots)$ 这样的符号，它给出第 $j$ 行中所有原子的自旋方向。我们假设只有相邻的层之间才有相互作用，即能量可以写成

\begin{equation*}
\mathcal{E} = \sum_{j=1}^{N} \mathcal{V}(\nu_j, \nu_{j+1}) + \sum_{j=1}^{N} \mathcal{V}(\nu_j)
\tag{10.88}
\end{equation*}

（其中包括每一层的“内能”，并且像 §1.6 中那样把两端闭合起来，使得 $\nu_{N+1} \equiv \nu_1$）。

于是系统的配分函数可以写成

\begin{equation*}
Z = \sum_{\nu_1}\sum_{\nu_2}\cdots\sum_{\nu_N}\left\{\prod_{j=1}^{N}\exp\left[-\left[\mathcal{V}(\nu_j, \nu_{j+1}) + \tfrac{1}{2}\mathcal{V}(\nu_j) + \tfrac{1}{2}\mathcal{V}(\nu_{j+1})\right]/kT\right]\right\}.
\tag{10.89}
\end{equation*}

但这不过相当于写下一串矩阵的乘积，其典型元素是

\begin{equation*}
P_{\nu_j\nu_{j+1}} = \exp\left\{-\left[\mathcal{V}(\nu_j, \nu_{j+1}) + \tfrac{1}{2}\mathcal{V}(\nu_j) + \tfrac{1}{2}\mathcal{V}(\nu_{j+1})\right]/kT\right\},
\tag{10.90}
\end{equation*}

这里要认识到，$\nu_j$ 是一个记号，按照层的复杂程度，它可以取两个或更多的值。实际上，式 (10.89) 说的无非是

\begin{align*}
Z &= \mathrm{trace}\,(P^N)\\
&= \sum_i \lambda_i^N,
\tag{10.91}
\end{align*}

其中 $\lambda_i$ 是式 (10.90) 中矩阵 $P$ 的本征值。实际上我们是假定 $N$ 很大——即我们的链几乎无限长——所以只有 $P$ 的最大本征值才是重要的。于是有结果

\begin{align*}
F &= -kT\ln Z\\
&\to -NkT\ln\lambda_{\max},
\tag{10.92}
\end{align*}

这表明自由能是一种广延性质。

现在，在一维链中，记号 $\nu_j$ 正是 $\sigma_j$ 的取值 $+1$ 和 $-1$。矩阵 (10.90) 可以写成

\begin{equation*}
P = \begin{pmatrix}\, y^{-\frac{1}{2}}z^{-\frac{1}{2}} & z^{\frac{1}{2}}\,\\[2pt]\, z^{\frac{1}{2}} & y^{\frac{1}{2}}z^{-\frac{1}{2}}\, \end{pmatrix}
\tag{10.93}
\end{equation*}

（用 (10.75) 的记号）。这个矩阵的本征值是

\begin{equation*}
\lambda_{\pm} = \tfrac{1}{2}\left\{(1+y) \pm \sqrt{\left[(1-y)^2 + 4yz^2\right]}\,\right\} y^{-\frac{1}{2}}z^{-\frac{1}{2}};
\tag{10.94}
\end{equation*}

把 $\lambda_+$ 代入 (10.92)，我们的问题便告解决。

这个结果的有趣之处在于，$F$ 既是 $y$ 的连续函数，又是 $z$ 的连续函数——也就是说，它是 $H$ 和 $T$ 的连续函数，没有奇点。零场下的比热可以算出，就是

\begin{equation*}
C_v = Nk\left(\frac{J}{kT}\right)^{2}\mathrm{sech}^2\left(\frac{J}{kT}\right),
\tag{10.95}
\end{equation*}

它在 $kT = J$ 附近有一个平滑的极大，但不表现出相变。

此外，在 $H = 0$ 时不存在与自发磁化强度相对应的解。\emph{线性链不是铁磁性的}。这一点从 Bragg--Williams 公式 (10.81) 本来看不出来，但在 Bethe 方法中确实显露出来：按式 (10.85)，当配位数 $p$ 趋于 2 时 $T_c \to 0$。当然，同样的结论对反铁磁性也成立。

\begin{figure}[H]
\centering
\includegraphics[width=0.72\textwidth]{fig_190.pdf}
\caption*{图 190\quad 线性链的比热。}
\end{figure}

重要的是要认识到这是一条\emph{拓扑}定理；Bethe 方法中 $p = 2$ 给不出转变温度这一点很可能只是偶然的。关键在于，线性链在任意一点都太容易被弄断。引入一个断口，便摧毁了长程序，使能量增加 $2J$；但断口可以出现在 $N$ 个格点中的任何一处，于是熵增加了 $k\ln N$。自由能的改变将会是

\begin{equation*}
F = 2J - kT\ln N,
\tag{10.96}
\end{equation*}

\begin{figure}[H]
\centering
\includegraphics[width=0.72\textwidth]{fig_191.pdf}
\caption*{图 191\quad 线性链中的长程序 (a) 被单个断口 (b) 破坏。}
\end{figure}

只要 $N$ 足够大，无论温度多低，它总可以被做成负的。由这个论证——或者由对一个像 (10.93) 那样的矩阵的本征值所作的更精细的讨论，该矩阵阶数有限而且所有元素都为正——可以推断：任何只在一个维度上无限的格子都将没有奇点。这是“一条链构不成晶体”这条格言的一个典型例证；凡研究凝聚态的学生都应当把它牢牢记在心头。

不难证明，在二维格子中上述反驳不再成立。例如，假设我们造出一个反转自旋的连通区域，如图 192 所示。这个区域将由一条长度为 $L$ 的多边形边界围住——也就是说，在两个区域之间的边界上将有 $L$ 条 AB 型的连键。于是系统的能量将增加 $2LJ$。但与这份能量相联系的熵是多少？一般地说，边界的每个节点处有 3 个可供选择的方向，因此铺排这条边界的方式约有 $3^L$ 种。实际上这是高估，因为多边形还必须重新闭合——但当 $L$ 很大时这个误差并不重要。于是，对自由能的贡献将是

\begin{align*}
F &\approx 2LJ - kT\ln 3^L\\
&= L(2J - kT\ln 3).
\tag{10.97}
\end{align*}

\begin{figure}[H]
\centering
\includegraphics[width=0.72\textwidth]{fig_192.pdf}
\caption*{图 192\quad 一个反转自旋的区域。}
\end{figure}

如果

\begin{equation*}
kT < \frac{2J}{\ln 3},
\tag{10.98}
\end{equation*}

则所有这类贡献都将为正。

有序态在这个温度以下应当是稳定的。

这个粗糙的论证可以由严格的计算来支持，因为 Ising 模型另一个可解的情形是二维正方格子。Onsager 证明了该系统在温度

\begin{equation*}
kT_c = \frac{2J}{\ln(1+\sqrt{2})}
\tag{10.99}
\end{equation*}

以下是“铁磁性的”（或“反铁磁性的”）；这一证明所要求的代数定理远远超出我们现在的范围（尽管其出发点实际上就是式 (10.92) 所引述的关系）。

这个结果最有趣的特征是：虽然自由能本身是温度的连续函数，比热却在居里温度两侧对数式地发散。这样的奇点实际上在任何固态体系中都没有观察到，但这也许是因为我们不容易真正实现一个二维的 Ising 自旋集合。

\begin{figure}[H]
\centering
\includegraphics[width=0.72\textwidth]{fig_193.pdf}
\caption*{图 193\quad 正方层格的比热：Onsager 公式。}
\end{figure}

Onsager 方法还可以应用于另外几种二维格子，它们表现出类似的性质；此外还有一些优美的拓扑定理，利用它们可以在其他情形下定出转变温度。然而，对于处在有限磁场中的二维 Ising 铁磁体，不存在严格解。

对三维格子来说，Ising 问题完全没有严格解。人们可以有把握地相信它们会表现出长程序，因为拓扑条件比二维更有利于合作效应；借助级数展开，居里点可以定得相当准确。

由这段讨论产生一条重要原理。按照 Bethe 方法，有序–无序转变应当只依赖于配位数 $p$（参看式 (10.75)）。这样一来，二维三角格子与三维简立方格子的配位数都是 $p = 6$，二者就应当有完全相同的性质。事实上，它们的行为大不相同：比热曲线不同，转变温度也不同。格子的维数对合作现象是带根本性的，实际上对固体的许多性质也是如此。

\begin{figure}[H]
\centering
\includegraphics[width=0.72\textwidth]{fig_194.pdf}
\caption*{图 194\quad 三角层格 (a) 与简立方格子 (b) 具有相同的配位数。}
\end{figure}

\section{自旋波}

让我们回到自旋哈密顿量 (10.29)，并假定交换积分为铁磁型的。在零磁场中

\begin{equation*}
\mathcal{H} = -\sum_{ll'} J_{ll'}\,\mathsf{S}_l \cdot \mathsf{S}_{l'}.
\tag{10.100}
\end{equation*}

设 $S$ 是每个离子的总自旋。我们可以用自旋的 $z$ 分量 $\mathsf{S}_l^z$ 的本征态来给第 $l$ 个离子上的自旋态分类。于是，对沿 $z$ 方向分量为 $S'$ 的态

\begin{equation*}
\mathsf{S}_l^z\left|S'\right\rangle_l = S'\left|S'\right\rangle_l
\tag{10.101}
\end{equation*}

（本节中取 $\hbar = 1$）。

铁磁性普遍理论的假定是：哈密顿量 $\mathcal{H}$ 有这样一个基态，其中所有自旋都沿 $z$ 方向最大限度地排列。我们假定这个态可以写成

\begin{equation*}
\left|0\right\rangle = \left|S\right\rangle_1\left|S\right\rangle_2\left|S\right\rangle_3\ldots\left|S\right\rangle_l\ldots\left|S\right\rangle_N,
\tag{10.102}
\end{equation*}

其中 $\mathsf{S}^z$ 在格子的每个格点上都取其最大值 $S$。

容易验证这是 $\mathcal{H}$ 的一个本征态。让我们为某一格点上的自旋定义两个新算符（暂时略去它的标记 $l$）

\begin{equation*}
\mathsf{S}^{+} \equiv \mathsf{S}^{x} + i\mathsf{S}^{y}, \quad \mathsf{S}^{-} \equiv \mathsf{S}^{x} - i\mathsf{S}^{y}.
\tag{10.103}
\end{equation*}

它们是自旋偏离（spin-deviation）算符。例如，考察一下用 $\mathsf{S}^{+}$ 作用到 $\mathsf{S}^z$ 的一个本征态上的效果；其结果是 $\mathsf{S}^z$ 的另一个本征态。于是，采用 (10.101) 的记号，

\begin{align*}
\mathsf{S}^z\{\mathsf{S}^{+}\left|S'\right\rangle\} &= \mathsf{S}^z(\mathsf{S}^{x}+i\mathsf{S}^{y})\left|S'\right\rangle\\
&= \{(\mathsf{S}^{x}\mathsf{S}^z+i\mathsf{S}^{y})+i(\mathsf{S}^{y}\mathsf{S}^z-i\mathsf{S}^{x})\}\left|S'\right\rangle\\
&= \mathsf{S}^{+}\mathsf{S}^z\left|S'\right\rangle+\mathsf{S}^{+}\left|S'\right\rangle\\
&= \mathsf{S}^{+}(S'+1)\left|S'\right\rangle\\
&= (S'+1)\{\mathsf{S}^{+}\left|S'\right\rangle\},
\tag{10.104}
\end{align*}

这里用了自旋算符的熟知对易关系。这样，$\mathsf{S}^{+}\left|S'\right\rangle$ 是 $\mathsf{S}^z$ 的本征值为 $S'+1$ 的本征态——换句话说，它就是态 $\left|S'+1\right\rangle$。$\mathsf{S}^{+}$ 的效果是使自旋在 $z$ 方向上的分量增加一个单位。类似地，$\mathsf{S}^{-}$ 使 $S'$ 减少一个量子。

\begin{figure}[H]
\centering
\includegraphics[width=0.72\textwidth]{fig_195.pdf}
\caption*{图 195\quad 算符 $\mathsf{S}_l^{\dagger}\mathsf{S}_{l'}^{-}$ 交换自旋偏离。}
\end{figure}

把 (10.103) 代入哈密顿量是很容易的事，后者可以写成

\begin{equation*}
\mathcal{H} = -\sum_{ll'} J_{ll'}\{\mathsf{S}_l^z\mathsf{S}_{l'}^z + \tfrac{1}{2}(\mathsf{S}_l^{+}\mathsf{S}_{l'}^{-}+\mathsf{S}_l^{-}\mathsf{S}_{l'}^{+})\}.
\tag{10.105}
\end{equation*}

当 $\mathcal{H}$ 作用在态 (10.102) 上时，只有 $z$ 分量作出贡献：

\begin{equation*}
\mathcal{H}\left|0\right\rangle = -\sum_{ll'} J_{ll'}S^2\left|0\right\rangle.
\tag{10.106}
\end{equation*}

算符 $\mathsf{S}_l^{+}$ 作用在态 $\left|S\right\rangle_l$ 上给出零，因为我们不能把 $\mathsf{S}_l^z$ 的本征值增大到超过它的最大值 $S$。于是，有序态 $\left|0\right\rangle$ 是总哈密顿量的本征态。

但现在来考虑系统的激发。自然而然的设想是，它们对应于单个自旋发生偏离的态，如图 195(a) 那样。然而，这样的态并不是 $\mathcal{H}$ 的本征态；算符 $\mathsf{S}_l^{+}\mathsf{S}_{l'}^{-}$ 作用在这个态上，会把它变成另一个态，如图 195(b) 那样，自旋偏离已经被交换到它的近邻身上。Ising 模型中所用的那种态的分类（例如图 185）在这里行不通；自旋偏离会在格子中不停地四处游走。

这种交换效应颇像 Bloch 态的紧束缚模型（§3.4）中电子从一个原子被递交给下一个原子的过程，而且结果证明它具有非常相似的数学性质。这一效应的数学表述如下。

让我们设 $S$ 是一个相当大的数（比如 $\frac{5}{2}$）。$\mathsf{S}^{-}$ 的效果是“产生一个自旋偏离”。我们不采用 (10.101) 的分类，而写

\begin{equation*}
n = S - S',
\tag{10.107}
\end{equation*}

并把态 $\left|n\right\rangle$ 定义为带有 $n$ 个自旋偏离量子的态。只要 $n < 2S$，我们就有通过 $\mathsf{S}^{-}$ 和 $\mathsf{S}^{+}$ 的作用来产生和湮没这些量子的自由。因此，这些算符带有简谐振子的常规理论中湮没算符和产生算符的某些性质（参看 §2.11）。

为了形式地表明这一点，我们可以计算对易子

\begin{align*}
[\mathsf{S}^{+},\mathsf{S}^{-}] &= i[\mathsf{S}^{y},\mathsf{S}^{x}] - i[\mathsf{S}^{x},\mathsf{S}^{y}]\\
&= -2i(i\mathsf{S}^z)\\
&= 2\mathsf{S}^z.
\tag{10.108}
\end{align*}

于是，若令

\begin{equation*}
a = \frac{\mathsf{S}^{+}}{(2\mathsf{S}^z)^{\frac{1}{2}}}, \quad a^{*} = \frac{\mathsf{S}^{-}}{(2\mathsf{S}^z)^{\frac{1}{2}}},
\tag{10.109}
\end{equation*}

则 $a$ 和 $a^{*}$ 具有标准的对易关系

\begin{equation*}
[a, a^{*}] = 1,
\tag{10.110}
\end{equation*}

而且互为复共轭。

定义 (10.109) 没有指明 $\mathsf{S}^{+}$ 和 $[\mathsf{S}^z]^{-\frac{1}{2}}$ 究竟应当按什么次序取——由于这些算符互不对易，这本来可能造成严重的后果。但若 $S$ 足够大，我们就可以忽略这个误差，把算符 $2\mathsf{S}^z$ 换成数 $2S$，因为我们只打算考虑几乎完全排列整齐的态。在这个近似下，自旋偏离态 $\left|n\right\rangle$ 是 $a^{*}a$ 的本征态，我们可以写

\begin{equation*}
a\left|n\right\rangle = n^{\frac{1}{2}}\left|n-1\right\rangle, \quad a^{*}\left|n\right\rangle = (n+1)^{\frac{1}{2}}\left|n+1\right\rangle.
\tag{10.111}
\end{equation*}

$n$ 的定义还意味着我们可以写

\begin{equation*}
\mathsf{S}^z = S - a^{*}a.
\tag{10.112}
\end{equation*}

现在把 (10.109) 和 (10.112) 代入哈密顿量 (10.105)。它近似地变成

\begin{align*}
\mathcal{H} &\approx -\sum_{ll'} J_{ll'}\{(S-a_l^{*}a_l)(S-a_{l'}^{*}a_{l'})+\tfrac{1}{2}\cdot 2S(a_l^{*}a_{l'}+a_l a_{l'}^{*})\}\\
&\approx -\sum_{ll'} J_{ll'}\{S^2 + S(a_l a_{l'}^{*}+a_l^{*}a_{l'}-a_l^{*}a_l-a_{l'}^{*}a_{l'})\},
\tag{10.113}
\end{align*}

这里我们略去了一个 $a_l^{*}a_l a_{l'}^{*}a_{l'}$ 量级的项。

式 (10.113) 中的第一项就是基态（式 (10.116)；译注：疑为 (10.106) 之误印）的能量；其余各项是算符，明确地表达自旋偏离从一个格点向下一个格点的传递。这颇像晶格动力学中的情形——晶格位移依靠弹性力从一个格点传到另一个格点——或者，正如我们已指出过的，像电子波的紧束缚公式。我们可以用在 §2.1 中用过的同一手法求出这个哈密顿量的本征态：代入

\begin{equation*}
a_l = N^{-\frac{1}{2}}\sum_{\mathbf{q}} a_{\mathbf{q}}\,e^{i\mathbf{q}\cdot\mathbf{l}}, \quad a_l^{*} = N^{-\frac{1}{2}}\sum_{\mathbf{q}} a_{\mathbf{q}}^{*}\,e^{-i\mathbf{q}\cdot\mathbf{l}},
\tag{10.114}
\end{equation*}

来代替每个格点上的湮没算符和产生算符。这等价于 (2.8)；波矢 $\mathbf{q}$ 满足与晶格动力学中相同的条件，而新算符 $a_{\mathbf{q}}$ 和 $a_{\mathbf{q}}^{*}$ 满足对易关系 (10.110)。

把 (10.114) 代入 (10.113)，并利用 $J_{ll'}$ 只依赖于 $\mathbf{h} = \mathbf{l}-\mathbf{l}'$ 这一事实（如 (2.6) 中那样），我们得到

\begin{equation*}
\mathcal{H} = \mathcal{H}_0 + \sum_{\mathbf{q}}\left\{\sum_{\mathbf{h}} 2SJ(\mathbf{h})\left(1-e^{-i\mathbf{q}\cdot\mathbf{h}}\right)\right\}a_{\mathbf{q}}^{*}a_{\mathbf{q}},
\tag{10.115}
\end{equation*}

其中 $\mathcal{H}_0$ 与算符 $a_{\mathbf{q}}$ 无关。但 $\mathcal{H}-\mathcal{H}_0$ 是一组简谐振子的哈密顿量，每个波矢 $\mathbf{q}$ 对应一个振子，每个振子具有能级

\begin{equation*}
\mathcal{E}_{\mathbf{q}} = n_{\mathbf{q}}\sum_{\mathbf{h}} 2SJ(\mathbf{h})\left(1-e^{-i\mathbf{q}\cdot\mathbf{h}}\right).
\tag{10.116}
\end{equation*}

换句话说，算符 $a_{\mathbf{q}}$、$a_{\mathbf{q}}^{*}$ 是自旋波的湮没算符和产生算符，自旋波是系统的元激发。第 $\mathbf{q}$ 个自旋波模中的能量以下列量子为单位而量子化：

\begin{equation*}
\hbar\omega_{\mathbf{q}} = \sum_{\mathbf{h}} 2SJ(\mathbf{h})\left(1-e^{-i\mathbf{q}\cdot\mathbf{h}}\right).
\tag{10.117}
\end{equation*}

仿照光子、声子和等离激元的命名，这样的元激发有时叫作磁振子（magnon）。

对于只含最近邻相互作用的简单格子，式 (10.117) 在小波矢下的形式容易算出：

\begin{equation*}
\hbar\omega_{\mathbf{q}} \approx 2SJq^2a^2,
\tag{10.118}
\end{equation*}

其中 $a$ 是晶格常数。于是，磁振子的能量正比于其波矢的平方——颇像电子的能量。实际上，可以定义一个“有效质量” $m^{*}$，使得

\begin{equation*}
\hbar\omega_{\mathbf{q}} \approx \frac{\hbar^2 q^2}{2m^{*}}.
\tag{10.119}
\end{equation*}

对于居里温度为几百开的典型铁磁体（由此可从式 (10.34) 得出 $J$ 的量级），我们发现 $m^{*}$ 约为自由电子质量的 10 倍。

激发谱的形式带来若干简单的推论。例如，由于磁振子服从 Bose--Einstein 统计，可以用 (2.46) 计算每个模的平均占据数：

\begin{equation*}
\overline{n}_{\mathbf{q}} = \frac{1}{e^{\hbar\omega_{\mathbf{q}}/kT}-1}.
\tag{10.120}
\end{equation*}

如果接受 (10.118)，并计算温度 $T$ 下所有被激发模中的磁振子总数，就得到

\begin{align*}
\overline{n} &= \frac{1}{8\pi^3}\int_0^{\infty}\frac{4\pi q^2\,dq}{\exp(2SJa^2q^2/kT)-1}\\
&= \frac{N}{4\pi^2}\left(\frac{kT}{2SJ}\right)^{\frac{3}{2}}\int_0^{\infty}\frac{x^{\frac{1}{2}}\,dx}{e^x-1}.
\tag{10.121}
\end{align*}

但每个磁振子都对应于一个自旋偏离的激发；它使系统的总自旋 $NS$ 减少一个单位（用 (10.112) 和 (10.114) 可以验证）。于是，自旋波诸模中 $\overline{n}$ 个量子的激发使总自旋减少 $\overline{n}$；磁化强度应当随温度的升高而按公式

\begin{equation*}
M(T) = M(0)\left\{1-\gamma\left(\frac{kT}{2SJ}\right)^{\frac{3}{2}}\right\},
\tag{10.122}
\end{equation*}

减小，其中 $M(0)$ 是 $T=0$ 时的饱和磁化强度，$\gamma$ 是一个可以由定积分 (10.121) 的值算出的数。

色散公式 (10.118) 的另一个简单推论是，磁振子应当对比热有贡献。显然，把 $\hbar\omega_{\mathbf{q}}$ 放到 (10.121) 的积分号下（参看式 (2.49)），就会在计算 $\mathcal{E}$ 时多引进 $T$ 的一次幂。这多出的一份 $T$ 的幂会在求微商得出比热时被取出来，因此比热也按 $T^{\frac{3}{2}}$ 变化。

铁磁介质中的自旋波还有许多其他有趣的性质。例如，它们可以与声子相互作用，并且可以像格波一样引起中子的非弹性散射。§2.8 的一般论证在此适用，并给出能量与晶体动量的选择定则。

从一般理论的观点看，我们在本节中所做的无非是对自旋哈密顿量 (10.100) 的态验证了 Bloch 定理（§1.4）。这个哈密顿量满足晶格平移不变性的条件，所以各态必定可以按波矢分类，并且必定具有满足 §2.5 各定律的谱。尽管如此，把哈密顿量约化成 (10.113) 的形式时不得不丢掉大量的项，特别是从自旋算符 $\mathsf{S}^{+}$ 和 $\mathsf{S}^{-}$ 定义算符 $a$ 和 $a^{*}$ 时所涉及的那些项。这些项给出自旋波之间的相互作用，在 (10.122) 按 $T$ 之幂的展开式中产生更多的项。事实是，在单个格点上能够产生的自旋偏离数以 $2S$ 为限。若产生高密度的磁振子，就有超越这一限度的趋势，可叠加性与独立性的假设便不再成立。实际上，有序态崩溃，我们穿过居里点；整个近似方案土崩瓦解。

我们对铁磁磁振子谱的推导是从 §10.4 的 Heisenberg 模型出发的。但是，物理上等价的现象也出现在 §10.5 的巡游电子模型中，尽管自旋并不局域在格点上。要构造行列式波函数 $\Psi(\mathbf{k},\mathbf{q})$——把一个电子从某个态 $\left|\mathbf{k},+\right\rangle$ 放进自旋已反转的另一个态 $\left|\mathbf{k}+\mathbf{q},-\right\rangle$——显然需要量级为交换能 $U$ 的能量。但电子气实际上是多体系统，除单粒子激发外还有集体模（参看 §5.7）。为了描述自旋波，我们把这些行列式作线性组合，

\begin{equation*}
\Psi_{\mathbf{q}} = \sum_{\mathbf{k}} A_{\mathbf{k}}\,\Psi(\mathbf{k},\mathbf{q})
\tag{10.123}
\end{equation*}

并选取系数 $A_{\mathbf{k}}$，使得 $\Psi_{\mathbf{q}}$ 是系统总哈密顿量的近似本征态。结果看上去与 (10.118) 完全一样；当 $\mathbf{q}\to 0$ 时能量趋于零。根据关于破缺对称效应的 \emph{Goldstone 定理}，这是其
必然后果：Heisenberg 哈密顿量 (10.29) 在磁场方向的转动下具有不变性。

从宏观上看，这一点是显而易见的：在长波长的自旋波中，介质在局域上被铁磁性地极化，相邻点之间相对自旋方向不会发生突然的反转。因此，磁振子是一种集体激发，其中电子的自旋在空间中强烈地相互关联。由于这个原因，磁振子–中子相互作用的一般理论并不依赖于电子是局域的还是巡游的（itinerant），因为衍射现象只依赖于这类双粒子的空间与时间关联（§2.8）。反铁磁性的自旋密度波（spin-density wave）理论（§10.7）中隐含的自旋关联，同样能为分析晶体中的磁有序图样给出正确的中子衍射结果（§2.6）。

\section{反铁磁基态}

初看之下，把 §10.11 的方法应用于 Heisenberg 反铁磁体似乎并不困难。假设我们有一个由 Ising 自旋构成的格子，它具有由变量 $\sigma_l$ 所定义的有序排列：在“自旋向上”的格点上 $\sigma_l$ 取 $+1$，在“自旋向下”的格点上取 $-1$。于是我们几乎可以完全仿照 (10.105) 写出我们的哈密顿量
\begin{equation*}
\mathcal{H} = -\sum_{ll'} J_{ll'} \bigl\{ \mathsf{S}_l^z \mathsf{S}_{l'}^z + \tfrac{1}{2}(\mathsf{S}_l^+ \mathsf{S}_{l'}^- + \mathsf{S}_l^- \mathsf{S}_{l'}^+) \bigr\} - \beta \sum_l (\sigma_l H_A \mathsf{S}_l^z + \mathbf{H}\cdot\mathsf{S}_l). \tag{10.124}
\end{equation*}

在这个表达式中，除了外磁场 $\mathbf{H}$ 之外，我们还纳入了一个“各向异性场”（anisotropy field）$H_A$，它总是平行于局域有序自旋的方向。于是，$H_A$ 在“向上”的格点上“向上”，在“向下”的格点上“向下”。这个场给出了量子化 $z$ 轴的方向（它不一定就是外加场的方向），除此之外它是任意的；我们很快就会看到需要它的理由。
\begin{figure}[H]
\centering
\includegraphics[width=0.72\textwidth]{fig_196.pdf}
\caption*{图 196\quad 反铁磁系统的量子化轴。}
\end{figure}

现在我们相对于对有序“Ising”态的偏离来定义局域的湮没算符与产生算符。在 $\sigma_l$ 为 $+1$ 的格点上，定义与 (10.109) 相同；而在“自旋向下”的格点上，我们要互换 $\mathsf{S}_l^+$ 与 $\mathsf{S}_l^-$ 的角色，因为我们是从第 $l$ 个离子的态 $|{-S}\rangle_l$ 出发来计量的。于是，在每一个格点上——无论“向上”还是“向下”——如下方程
\begin{align*}
\left.
\begin{aligned}
\mathsf{S}_l^x &\approx (\tfrac{1}{2}S)^{\frac{1}{2}}(a_l + a_l^*), \\
\mathsf{S}_l^y &\approx -i\sigma_l (\tfrac{1}{2}S)^{\frac{1}{2}}(a_l - a_l^*), \\
\mathsf{S}_l^z &= \sigma_l (S - a_l^* a_l),
\end{aligned}
\right\} \tag{10.125}
\end{align*}
定义了湮没和产生自旋偏离的算符 $a_l$ 与 $a_l^*$。（在某些论述中，人们在两个不同的子晶格（sublattice）上定义不同的算符；不过由于当 $l \neq l'$ 时 $a_l$ 与 $a_l^*$ 对易，这只是标签的问题。）

把 (10.125) 代入 (10.124) 是相当容易的。第一项就是这个态的“Ising”能量
\begin{equation*}
\mathcal{E}_0 = -\sum_{ll'} J_{ll'}\, \sigma_l \sigma_{l'} S^2 - N\beta S H_A \tag{10.126}
\end{equation*}
这一项我们可以丢弃。其余的项——在 $a_l$、$a_l^*$ 中保留到二阶乘积——虽然复杂，但很容易写出。

仿照 (10.114) 的一个 Fourier 变换引入新的湮没与产生算符 $b_{\mathbf{q}}$ 和 $b_{\mathbf{q}}^*$，使得
\begin{equation*}
a_l = N^{-\frac{1}{2}} \sum_{\mathbf{q}} b_{\mathbf{q}}\, e^{i\mathbf{q}\cdot\mathbf{l}}, \qquad a_l^* = N^{-\frac{1}{2}} \sum_{\mathbf{q}} b_{\mathbf{q}}^*\, e^{-i\mathbf{q}\cdot\mathbf{l}}. \tag{10.127}
\end{equation*}
除了“Ising”能量 (10.126) 之外，我们发现哈密顿量可以写成
\begin{equation*}
\mathcal{H} = \tfrac{1}{2} \sum_{\mathbf{q}} \bigl\{ A_{\mathbf{q}} (b_{\mathbf{q}}^* b_{\mathbf{q}} + b_{-\mathbf{q}}^* b_{-\mathbf{q}}) + B_{\mathbf{q}} (b_{\mathbf{q}} b_{-\mathbf{q}} + b_{\mathbf{q}}^* b_{-\mathbf{q}}^*) \bigr\} + \mathbf{H}\cdot\mathbf{M}, \tag{10.128}
\end{equation*}
其中系数与 (2.12) 或 (10.116) 相类似：——
\begin{align*}
\left.
\begin{aligned}
A_{\mathbf{q}} &= \beta H_A + \sum_{\mathbf{h}} 2S J(\mathbf{h}) \bigl\{ \sigma_{\mathbf{h}} - \tfrac{1}{2}(1 + \sigma_{\mathbf{h}})\cos\mathbf{q}\cdot\mathbf{h} \bigr\}, \\
B_{\mathbf{q}} &= -\sum_{\mathbf{h}} S J(\mathbf{h})(1 - \sigma_{\mathbf{h}})\cos\mathbf{q}\cdot\mathbf{h}.
\end{aligned}
\right\} \tag{10.129}
\end{align*}
在这些表达式中，$\sigma_{\mathbf{h}}$ 表示位于 $\mathbf{h}$ 处的格点相对于中心格点符号的符号。我们可以把 $J(\mathbf{h})$ 的力程比最近邻相互作用更长、甚至在某些距离上为铁磁型的情形也包括进来；唯一必要的条件是 $\sigma_l$ 应当正确地描述 Ising 模型的基态。

在 (10.128) 中，外磁场与一个复杂的算符 $\mathbf{M}$ 联结在一起出现；我们不打算一般地定义它，尽管在简单的情形中可以看出，它对应于一个在两个子晶格的格点上（沿 $z$ 轴）看似正负交替的矩。

这是我们在 (10.125) 中对局域量子化轴的任意选取所造成的结果——对这个项我们就不再深究了。

我们的“交替”态的另一个效应是：(10.128) 中存在把 $b_{\mathbf{q}}$ 与 $b_{-\mathbf{q}}$ 联结起来的项。当我们把 (10.125) 代入 (10.124) 时，它们通过像 $a_l a_{l'}$ 这样的乘积产生，而这类乘积在铁磁态 (10.113) 中并不出现。强加上一个 Ising 自旋背景系统降低了体系的基本平移对称性，因此哈密顿量不再能被 Fourier 变换 (10.127) 自动对角化。或许我们应当取每个晶胞两个自旋，并作一种颇不相同的 Fourier 分析。

然而，这个问题实际上并不比晶格动力学的类似问题中相互作用矩阵 (2.13) 的约化更困难。我们通过一个正则变换引入新变量 $c_{\mathbf{q}}$、$c_{\mathbf{q}}^*$：
\begin{align*}
\left.
\begin{aligned}
b_{\mathbf{q}} &= c_{\mathbf{q}}\cosh\theta_{\mathbf{q}} + c_{-\mathbf{q}}^*\sinh\theta_{\mathbf{q}}, \qquad b_{\mathbf{q}}^* = c_{\mathbf{q}}^*\cosh\theta_{\mathbf{q}} + c_{-\mathbf{q}}\sinh\theta_{\mathbf{q}}, \\
b_{-\mathbf{q}} &= c_{\mathbf{q}}^*\sinh\theta_{\mathbf{q}} + c_{-\mathbf{q}}\cosh\theta_{\mathbf{q}}, \qquad b_{-\mathbf{q}}^* = c_{\mathbf{q}}\sinh\theta_{\mathbf{q}} + c_{-\mathbf{q}}^*\cosh\theta_{\mathbf{q}},
\end{aligned}
\right\} \tag{10.130}
\end{align*}
它保持对易子
\begin{equation*}
[c_{\mathbf{q}}, c_{\mathbf{q}}^*] = \cosh^2\theta_{\mathbf{q}} - \sinh^2\theta_{\mathbf{q}} = 1. \tag{10.131}
\end{equation*}
参数 $\theta_{\mathbf{q}}$ 表示在这四个算符的空间中绕一个虚角度的转动。值得注意的是，$c_{\mathbf{q}}$ 的共轭是 $c_{-\mathbf{q}}^*$；我们必须反转运动方向，才能使它看上去像是给定波的物理共轭。

如果把 (10.130) 代入 (10.128)，就会得到一组“非对角”项
\begin{equation*}
\tfrac{1}{2} \sum_{\mathbf{q}} \bigl\{ 2A_{\mathbf{q}}\cosh\theta_{\mathbf{q}}\sinh\theta_{\mathbf{q}} + B_{\mathbf{q}}(\cosh^2\theta_{\mathbf{q}} + \sinh^2\theta_{\mathbf{q}}) \bigr\} (c_{\mathbf{q}} c_{-\mathbf{q}} + c_{\mathbf{q}}^* c_{-\mathbf{q}}^*), \tag{10.132}
\end{equation*}
只要把 $\theta_{\mathbf{q}}$ 选得满足
\begin{align*}
\frac{B_{\mathbf{q}}}{A_{\mathbf{q}}} &= -\frac{2\cosh\theta_{\mathbf{q}}\sinh\theta_{\mathbf{q}}}{\cosh^2\theta_{\mathbf{q}} + \sinh^2\theta_{\mathbf{q}}} \\
&= -\tanh 2\theta_{\mathbf{q}}, \tag{10.133}
\end{align*}
这组项便可以对每一个 $\mathbf{q}$ 值都变为零。对角项于是给出如下结果
\begin{equation*}
\mathcal{H} = \sum_{\mathbf{q}} \bigl[ (A_{\mathbf{q}}^2 - B_{\mathbf{q}}^2)^{\frac{1}{2}} (c_{\mathbf{q}}^* c_{\mathbf{q}} + \tfrac{1}{2}) - A_{\mathbf{q}} \bigr], \tag{10.134}
\end{equation*}
其中我们已利用对易关系把算符排成了这一顺序。现在的求和遍及 $\mathbf{q}$ 的一切取值，正负兼备。

这个结果与 (10.115) 十分相似。$c_{\mathbf{q}}^* c_{\mathbf{q}}$ 的本征态，比方说记作 $|n_{\mathbf{q}}\rangle$，对应于该模中激发了 $n_{\mathbf{q}}$ 个量子。在简单情形下，反铁磁振子（antiferromagnon）的能量将是
\begin{align*}
\hbar\omega_{\mathbf{q}} &= (A_{\mathbf{q}}^2 - B_{\mathbf{q}}^2)^{\frac{1}{2}} \\
&= \bigl\{ (\beta H_A - 2zSJ)^2 - (2pSJ\langle\cos qa\rangle)^2 \bigr\}^{\frac{1}{2}}, \tag{10.135}
\end{align*}
其中 $\langle\cos qa\rangle$ 是这个量对 $p$ 个最近邻取的平均值，这些近邻与原点相距 $a$，我们假定它们通过交换积分 $J$ 相互作用。如果写成
\begin{equation*}
\hbar\omega_A = \beta H_A, \quad \hbar\omega_J = -2pSJ \tag{10.136}
\end{equation*}
（记得 $J$ 在反铁磁性中是负的），则
\begin{equation*}
\omega_{\mathbf{q}} = \bigl\{ (\omega_A + \omega_J)^2 - (\omega_J\langle\cos qa\rangle)^2 \bigr\}^{\frac{1}{2}} \tag{10.137}
\end{equation*}
就是自旋波谱。

这看上去相当简单。例如，设 $\omega_A \ll \omega_J$，那么
\begin{align*}
\omega_{\mathbf{q}} &\approx \omega_J\bigl\{ 1 - \langle\cos qa\rangle^2 \bigr\}^{\frac{1}{2}} \\
&\approx \frac{1}{\sqrt{3}}\,\omega_J\, .\, qa, \tag{10.138}
\end{align*}
对应于频率对波数呈线性依赖，正如晶格谱 (2.21) 中那样。磁振子对比热的贡献应当按 $T^3$ 变化，也正如 (2.57) 中那样。

让我们试着遵循 (10.122)，计算每个子晶格（sublattice）的磁化强度随温度的变化。平均自旋偏离显然是
\begin{align*}
\langle a_l^* a_l\rangle &= \frac{1}{N}\sum_{\mathbf{q}} \langle b_{\mathbf{q}}^* b_{\mathbf{q}}\rangle \\
&= \frac{1}{N}\sum_{\mathbf{q}} \bigl\langle \cosh^2\theta_{\mathbf{q}}\, c_{\mathbf{q}}^* c_{\mathbf{q}} + \sinh^2\theta_{\mathbf{q}}\, c_{\mathbf{q}} c_{\mathbf{q}}^* \bigr\rangle \\
&= \frac{1}{N}\sum_{\mathbf{q}} \Bigl\{ \bigl(n_{\mathbf{q}} + \tfrac{1}{2}\bigr)\frac{A_{\mathbf{q}}}{\hbar\omega_{\mathbf{q}}} - \tfrac{1}{2} \Bigr\}, \tag{10.139}
\end{align*}
这里用到了变换 (10.127) 与 (10.130)。在此式中 $n_{\mathbf{q}}$ 是第 $q$ 个模中激发的磁振子数；我们自然要用 Bose–Einstein 分布 (10.120) 把它作为温度的函数计算出来。

但是即使在绝对零度，平均自旋偏离也不会消失。在 (10.139) 中取 $n_{\mathbf{q}} = 0$，我们仍然有
\begin{align*}
\langle a_l^* a_l\rangle_{T=0} &= \frac{1}{2N}\sum_{\mathbf{q}} \left( \frac{A_{\mathbf{q}}}{\hbar\omega_{\mathbf{q}}} - 1 \right) \\
&\approx \frac{1}{2N}\sum_{\mathbf{q}} \Bigl\{ \frac{1}{\bigl(1 - \langle\cos qa\rangle^2\bigr)^{\frac{1}{2}}} - 1 \Bigr\} \tag{10.140}
\end{align*}
（在 (10.138) 所描述的简单情形下）。这个和是可以计算的；它的行为如同积分
\begin{equation*}
\int \frac{1}{q}\,\mathrm{d}\mathbf{q} \tag{10.141}
\end{equation*}
只有当 $\mathrm{d}\mathbf{q}$ 是在二维或三维流形上取值时，它才收敛。于是我们证实：一维反铁磁体即使在 $T = 0$ 时也不具有稳定的有序态。这超出了 (10.96) 的结果——后者表明一维 Ising 模型对热涨落是不稳定的；在反铁磁情形，与自旋交换相联系的涨落会摧毁体系所假想的有序基态。

\begin{figure}[H]
\centering
\includegraphics[width=0.72\textwidth]{fig_197.pdf}
\caption*{图 197\quad 在反铁磁有序态 (a) 中，交换算符 $\mathsf{S}_l^+ \mathsf{S}_{l'}^-$ 产生两个自旋偏离 (b)。}
\end{figure}

即使在二维和三维情形，每个子晶格的磁矩也达不到其假定的极大值。事实上，由 (10.125) 中的 $\sigma_l$ 所定义的“交替”态并非真正是哈密顿量的本征态。我们可以试着用 (10.106) 的方法来核查这一点；正如在那里指出的，$\mathsf{S}_l^+ \mathsf{S}_{l'}^-$ 项倾向于把一个自旋偏离从一个格点传递到相邻的下一个格点。而在反铁磁情形，它会在两个相邻格点上同时产生自旋偏离——于是基态必定包含由这种组态及类似组态出发所能得到的一切状态的混合。

三维体系的反铁磁基态并非完全有序（当 $S = \tfrac{1}{2}$ 时，每个子晶格大概只能达到其极大磁矩的 90\%），这一事实也反映在能量的期望值上。从 (10.126)、(10.129) 与 (10.134) 可以看到，它是
\begin{align*}
\langle\mathcal{H}\rangle + \mathcal{E}_0 &= \sum_{\mathbf{q}} \bigl( n_{\mathbf{q}} + \tfrac{1}{2} \bigr) \hbar\omega_{\mathbf{q}} - A_{\mathbf{q}} + \mathcal{E}_0 \\
&= -2NpS(S+1)J + \sum_{\mathbf{q}} \bigl( n_{\mathbf{q}} + \tfrac{1}{2} \bigr) \hbar\omega_{\mathbf{q}}, \tag{10.142}
\end{align*}
即便在 $T = 0$、$n_{\mathbf{q}} = 0$ 时，它也大于人们对“Ising”态算出的对角值。磁振子各模中的“零点能”必须包括进去；它使基态的能量升高，基态能量位于 $-2NpJS(S+1)$ 与 $-2NpJS^2$ 之间的某处。

\begin{figure}[H]
\centering
\includegraphics[width=0.72\textwidth]{fig_198.pdf}
\caption*{图 198\quad 在没有各向异性场时，反铁磁阵列 (a) 的自旋可以在能量不变的情况下转动 (b)。}
\end{figure}

关于 (10.139) 还有进一步的一点值得注意。按照 (10.138)，$\omega_{\mathbf{q}}$ 在 $q = 0$ 处为零，这样，在平均自旋偏离的计算中这一项的贡献就会是无穷大。我们需要各向异性场（anisotropy field）$H_A$ 来使得
\begin{equation*}
\omega_{\mathbf{q}} \to \sqrt{\bigl( 2\omega_A\omega_J \bigr)} \quad \text{当}\; q \to 0 \tag{10.143}
\end{equation*}
从而避免这个奇点。事实是：我们的体系在其总自旋为零的基态下，对于如图 198 所示的所有自旋的整体均匀转动是不稳定的。各向异性场正是用来稳定量子化方向的。这个场可以非常小——没有 $H_A$，(10.141) 中的极限在 $q = 0$ 处也可以是有限的——但原则上各向异性场必须存在，而且实际上它也确实存在。

实际上，我们可以把频率 (10.143) 作为体系的反铁磁共振（antiferromagnetic resonance）频率观测到。即使 $H_A$ 只有几百高斯，这也是一个很高的频率，但它可以通过施加一个强的外磁场而被带入共振。

反铁磁有序态的磁化率规则不难核查；当然，关于磁振子–声子相互作用等等也已有各种计算。然而，最有意思的一点或许是：如果我们用 (10.139) 去计算有限温度下子晶格的极化，就会发现一个行为如下的表达式
\begin{align*}
\sum_{\mathbf{q}} \frac{\bar{n}_{\mathbf{q}} A_{\mathbf{q}}}{\hbar\omega_{\mathbf{q}}} &\approx \sum_{\mathbf{q}} \frac{kTA_{\mathbf{q}}}{(\hbar\omega_{\mathbf{q}})^2} \\
&\propto \int \frac{1}{q^2}\,\mathrm{d}\mathbf{q} \tag{10.144}
\end{align*}
它在二维情形下于 $q = 0$ 处对数发散。

这样，交换效应在受到热激发时，就足以摧毁二维 Ising 模型的有序反铁磁态；Onsager 结果在这个情形下对真实自旋并不成立（尽管它对铁磁性确实成立）。积分 (10.144) 在三维情形确实收敛，因此反铁磁态在三维下很可能是稳定的——总会有一个由 (10.35) 那样的偶极–偶极项产生的微小各向异性场。但这一点尚未得到严格证明；其正确性的最好明证，也许当属 §2.6 中提到过的中子衍射实验——它们在物理上直接显示出了两个磁子晶格。
